\documentclass[12pt]{article}
\usepackage[utf8]{inputenc}

\usepackage{amsmath,amssymb,bm}
\usepackage{paralist}
\usepackage{graphicx} %% add this and next lines if pictures should be in eps format
\usepackage{epsfig} %For pictures: screened artwork should be set up with an 85 or 100 line screen
%\usepackage{subfigure,appendix}
%\usepackage[dvips]{color,graphics}
\usepackage{psfrag}
\usepackage{multicol}
\usepackage{setspace,tabularx,fancyhdr}
\usepackage{ulem}
\usepackage[sort,compress]{natbib}
\usepackage{tocbibind}
\usepackage{color}
\usepackage[usenames,dvipsnames,svgnames,table]{xcolor}
\pagestyle{plain}
\usepackage[top=2cm,bottom=2cm,left=2cm,right=2cm]{geometry}
\usepackage{enumitem}
\setlist[itemize]{leftmargin=*}
\usepackage{wrapfig}
\usepackage{type1cm}
\usepackage{booktabs}
\usepackage{amsmath}
\usepackage{siunitx}
\usepackage{float}
\usepackage{subcaption}
\usepackage{url}

\usepackage{multirow}
\usepackage{booktabs}

\newlength\normalFontSize\newlength\fontSizeIncrement
\setlength\normalFontSize{12pt}

\def\beq{\begin{equation}}
\def\eeq{\end{equation}}
\def\bea{\begin{eqnarray}}
\def\eea{\end{eqnarray}}
\def\ergo{\Rightarrow}

\newcommand{\red}[1]{{\color{red} #1}}
\newcommand{\green}[1]{{\color{green} #1}}
\newcommand{\blue}[1]{{\color{blue} #1}}
\newcommand{\violet}[1]{{\color{violet} #1}}

\definecolor{myyellow}{rgb}{0.94, 0.86, 0.51}
\newcommand{\myyellow}[1]{{\color{myyellow} #1}}

\definecolor{mygreen}{rgb}{0.2, 0.8, 0.2}
\newcommand{\mygreen}[1]{{\color{mygreen} #1}}

\definecolor{mypink}{rgb}{0.99, 0, 0.99}
\newcommand{\mypink}[1]{{\color{mypink} #1}}

\definecolor{mypurple}{rgb}{0.75, 0, 0.75}
\newcommand{\mypurple}[1]{{\color{mypurple} #1}}

\definecolor{cadmiumorange}{rgb}{0.93, 0.53, 0.18}
\newcommand{\orange}[1]{{\color{cadmiumorange} #1}}

\newcommand{\cosmic}{\mypurple{c}\blue{o}\mygreen{s}\myyellow{m}\orange{i}\red{c}\,}

\newcommand{\crucial}{\mypurple{c}\blue{r}\mygreen{u}\myyellow{c}\orange{i}\red{a}\mypink{l}\,}

\newcommand{\paramount}{p\mypurple{a}\blue{r}\mygreen{a}\myyellow{m}\orange{o}\red{u}\mypink{n}t\,}

\newcommand{\transformative}{\mypurple{tr}\blue{an}\mygreen{sf}\myyellow{or}\orange{ma}\red{ti}\mypink{ve}\,}


\newcommand{\critical}{c\mypurple{r}\blue{i}\mygreen{t}\myyellow{i}\orange{c}\red{a}\mypink{l}\,}


\newcommand{\rep}[2]{\sout{#1}\red{#2}}

\newcommand{\FIX}[1]{{\color{red}$\blacksquare$~\textsf{[FIX: #1]}}}

\newcommand{\TODO}[1]{{\color{red}$\blacksquare$~\textsf{[TODO: #1]}}}

\newcommand{\CITE}{{\color{red}$\blacksquare$~\textsf{[CITATIONS NEEDED]}}}

\newcommand{\GOOD}[1]{{\color{blue}$\blacksquare$~\textsf{[#1]}}}

\newcommand{\IGNORE}[1]{}

\title{\red{BIU MSc Proposal Template \\ Note title page should be made in Word or something\\
in both Hebrew and English \\
following BIU template\\
with a BIU logo and all that, see:\\
\url{https://graduate-school.biu.ac.il/node/1772} \\
so the actual stuff starts from the next page}
}

\author{Dolev Azrad\\
Supervisor: Dr. Ofek Birnholtz\\
Department of Physics, Bar-Ilan University}

\begin{document}
\setcitestyle{square}

\maketitle

\clearpage

\tableofcontents

\clearpage



\section{Introduction}
Gravitational waves (GWs), predicted by Einstein's general theory of relativity \cite{Einstein1916} and first directly detected by the LIGO--Virgo collaboration in 2015 \cite{B.abbot et al2016, LIGOScientific2015, Acernese2015}, have opened a fundamentally new observational window on the universe. Unlike electromagnetic astronomy, gravitational-wave observations probe the dynamics of spacetime itself, granting direct access to the most extreme astrophysical environments: the final inspiral, merger, and ringdown of compact binary systems composed of black holes and neutron stars. With 90 gravitational-wave candidates catalogued in GWTC-3 \cite{GWTC3}, including 35 new detections from the second half of the third observing run alone, the field has transitioned from the era of individual landmark events to that of statistical population analysis.

Extracting astrophysical information from a gravitational-wave signal requires \textit{parameter estimation} (PE): the Bayesian inference of up to 15 source parameters — component masses, spin vectors, sky location, luminosity distance, inclination, orbital phase, and coalescence time — from detector strain data \cite{Romero-Shaw2020, Biscoveanu2020}. This is accomplished using stochastic samplers such as nested sampling or Markov Chain Monte Carlo (MCMC), which map the full 15-dimensional posterior probability distribution. For a single event, a production-level PE run with a precessing waveform model such as \texttt{IMRPhenomXPHM} requires 20--30 hours of continuous computation on a modern CPU \cite{Romero-Shaw2020}. Applied to a catalog of $\sim$200 events per observing run, this represents upwards of 6{,}000 CPU-hours of analysis per waveform family.

A second structural feature of LIGO–Virgo–KAGRA operations compounds this burden. Following each detection, the collaborations operate a tiered analysis pipeline. Within seconds of a candidate trigger, BAYESTAR produces a rapid sky-localization map; within tens of minutes, an automated low-latency PE run using \texttt{Bilby} together with the reduced-order quadrature (ROQ) likelihood produces a first set of posterior samples on the \textit{preliminary} power spectral density (PSD) \cite{Chaudhary2024, Bilby: A user-friendly, Smith2016}. These low-latency products drive multi-messenger alerts but are not the final, catalog-quality results. Once detector calibration is finalized — a process that revises the PSD amplitude by typically less than a few percent and the phase by less than a few degrees across the most sensitive band \cite{Sun2020, Biscoveanu2020} — a second, full-fidelity PE run is launched on the corrected PSD. This second run is what populates the GWTC catalogs and supports population-level inference. Under the current paradigm, however, the final run discards the low-latency posterior and re-samples the full astrophysical prior from scratch, duplicating the work already performed in the low-latency tier. For a catalog of $\sim$200 events per observing run, this duplication represents thousands of CPU-hours of avoidable cost.

This thesis proposes and validates a \textbf{two-phase Scout and Refine framework} that eliminates this duplication. In \textit{Phase 1} (Scout), a nested sampler with uninformative priors operates on the preliminary PSD — analogously to the existing low-latency tier — and identifies the high-probability region of parameter space at reduced computational cost. In \textit{Phase 2} (Refine), truncated Gaussian priors constructed from the Phase 1 posteriors constrain the final, full-fidelity sampler running on the calibration-finalized PSD. Because the Refine sampler explores a prior volume that is exponentially smaller than the full astrophysical prior, the wall-clock cost of producing catalog-quality posteriors is substantially reduced. The framework targets specifically the bottleneck of the final, catalog re-analysis run; it is not a replacement for the existing low-latency alert tier, which is already fast.

The central result of this work is a systematic validation of this framework across ten luminosity distances spanning three orders of magnitude in signal strength, from 10 Mpc (extreme SNR $\approx$ 790) to 5000 Mpc (sub-detection, SNR $\approx$ 2). We demonstrate that the two-phase method achieves a mean wall-clock reduction of \textbf{44\%} in the standard detection regime (40--150 Mpc) while maintaining mathematically unbiased Bayes Factors ($\Delta \ln \mathrm{BF} < 1$ everywhere above 40 Mpc) and statistically indistinguishable posteriors ($\Delta_{\mathrm{param}} < 1.0$ for all parameters at all distances). This thesis is organized as follows: Section~2 reviews the relevant background in GW PE and nested sampling; Section~3 states the research objectives; Sections~4--5 detail the hypotheses and methodology; Section~6 presents the results; Section~7 provides theoretical interpretation and discusses limitations; and Section~8 concludes.




\begin{figure}[ht]
\centering
\begin{subfigure}[t]{0.48\textwidth}
    \centering
    \includegraphics[width=\linewidth]{visualising_the_results_11_0.png}

    \label{fig:original_gw_recovery}
\end{subfigure}
\hfill
\begin{subfigure}[t]{0.48\textwidth}
    \centering
    \includegraphics[width=\linewidth]{visualise_results_final_corner.png}
    \label{fig:my_gw_recovery}
\end{subfigure}
\caption{Comparison of PE for a binary black hole gravitational-wave signal. Left: Results obtained using Bilby as reported in literature. Right: Results obtained by the author, replicating the analysis. On the left: Injecting and recovering a binary black hole gravitational-wave signal with Bilby. The result of PE, The posterior distributions are shown in blue and the injected values in orange.
    This figure is taken from Bilby: A user-friendly Bayesian inference library for gravitational-wave astronomy \cite{Bilby: A user-friendly}
On the right: This figure was created by me to duplicate the results.}
\label{fig:gw_recovery_comparison}
\end{figure}




\newpage
\section{Background \& Literature Review}

\subsection{Gravitational-Wave Parameter Estimation: Key Parameters and Their Significance}

In GW astronomy, accurately determining both intrinsic and extrinsic parameters is crucial for the analysis and interpretation of GW signals. Intrinsic parameters, such as the masses (\(M\)) and spins (\(\vec{S}\)) of celestial bodies, define the waveform's characteristics independently of the observer's perspective. Conversely, extrinsic parameters, including the luminosity distance (\(d_L\)) and sky location (\(\theta, \phi\)), influence the signal's perceived properties and detectability. The orientation of the binary system's orbital plane relative to the observer's line of sight, an extrinsic factor, alters the observed amplitude and polarization of the GW signal. Distinguishing between intrinsic qualities, intrinsic to the astrophysical sources, and extrinsic qualities, dependent on observational context, is essential for a thorough analysis of gravitational waves.

 \begin{figure}[ht]
\centering
\includegraphics[width=0.5\textwidth]{signal_Waveform.png}
\caption{On top: Estimated gravitational wave amplitude of GW150914 at the Hanford detector. Above that are the Schwarzschild horizons of both merging black holes shown as calculated numerically from the general theory of relativity.[Image: LIGO / Redesign: Daniela Leitner]\cite{Observation of Gravitational Waves from a Binary Black Hole Merger}}
\label{fig:signal_waveform}
    \end{figure}
\newpage


\textbf{\textit{ The parameters commonly estimated include:}}

\begin{table}[ht]
\centering
\begin{tabular}{@{} ll ll ll ll @{}}
\toprule
& Variable & Unit & Prior & BBH Range & NSBH Range & BNS Range \\
\midrule
\multirow{4}{*}{Intrinsic}
& \( m_{1,2} \) & \( M_{\odot} \) & Uniform & \( [5, 100] \) & \( [1.2, 10] / [5, 100] \) & \( [1.2, 2.5] \) \\
& \( a_{1,2} \) & -- & Uniform & \( [0, 0.98] \) & \( [0, 0.98] / [0, 0.05] \) & \( [0, 0.05] \) \\
& \( \theta_{1,2} \) & rad. & Sin & \( [0, \pi] \) & \( [0, \pi] \) & \( [0, \pi] \) \\
& \( \phi_{JL} \) & rad. & Uniform & \( [0, 2\pi] \) & \( [0, 2\pi] \) & \( [0, 2\pi] \) \\
\midrule % Visual separator between intrinsic and extrinsic parameters
\multirow{5}{*}{Extrinsic}
& \( d_L \) & Mpc & Comoving & \( [10^2, 10^4] \) & \( [10^2, 10^4] \) & \( [10^2, 10^4] \) \\
& \( ra \) & rad. & Uniform & \( [0, 2\pi] \) & \( [0, 2\pi] \) & \( [0, 2\pi] \) \\
& \( dec \) & rad. & Cos & \( [-\frac{\pi}{2}, \frac{\pi}{2}] \) & \( [-\frac{\pi}{2}, \frac{\pi}{2}] \) & \( [-\frac{\pi}{2}, \frac{\pi}{2}] \) \\
& \( \psi \) & rad. & Sin & \( [0, \pi] \) & \( [0, \pi] \) & \( [0, \pi] \) \\
& \( \phi_c \) & rad. & Uniform & \( [0, 2\pi] \) & \( [0, 2\pi] \) & \( [0, 2\pi] \) \\
\bottomrule
\end{tabular}
\caption[Priors for PE across GW source estimation]{Priors for PE across binary black hole (BBH), neutron star-black hole (NSBH), and binary neutron star (BNS) mergers. The table is organized to present both intrinsic (Top) and extrinsic (Bottom) parameters for each source type. Intrinsic variables include the masses \( m_{1,2} \), dimensionless spin magnitudes \( a_{1,2} \), tilt angles \( \theta_{1,2} \), and spin vectors \( \delta\phi_{1,2} \) and \( \delta\phi_{JL} \). Extrinsic parameters encompass the luminosity distance \( d_L \), right ascension \( ra \), declination \( dec \), inclination angle \( \iota \), polarization angle \( \psi \), and phase at coalescence \( \phi_c \), with all phases, spins, and inclination angles defined at a reference frequency. The 'Sin' and 'Cos' priors correspond to distributions uniform in sine and cosine, while 'comoving' denotes a uniform distribution in comoving volume. This table extends the framework provided by \cite{Bilby: A user-friendly} to encompass the parameter space relevant for all three types of gravitational wave sources, reflecting the diverse nature of these cosmic events. See figure \ref{fig:params4} for visual representation.}
\label{table:gw_parameters_ranges}
\end{table}
 \begin{figure}[ht]
\centering
\includegraphics[width=0.99\textwidth]{params4.png}

\caption[Schematic of a binary black hole system]{
Illustration of a binary black hole (BBH) system depicting key dynamical and spin parameters. Intrinsic parameters shown include the black hole masses (\( m_{1} \) and \( m_{2} \)), spin vectors (\( \vec{S}_{1} \) and \( \vec{S}_{2} \)), and their respective orientations (\( \theta_{1,2} \) relative to the orbital angular momentum, not explicitly labeled in the image but implied by the spin vector orientation). Extrinsic parameters like the system's position and orientation in space, though not labeled, can be inferred from the relative positions and momenta (\( \vec{P}_{1} \) and \( \vec{P}_{2} \)), and the center of mass (CM). This schematic provides a visual representation of the parameters important for the estimation process in gravitational wave astronomy as listed in the accompanying PE table. Note that while the table includes ranges for NSBH and BNS systems, this illustration focuses on BBH parameters, which are directly comparable to those listed for BBH in the table. figure courtesy of Ofek Dan, BIU
}
\label{fig:params4}
\end{figure}

\begin{itemize}


    \item \textbf{Masses of the Objects}: In gravitational wave astronomy, understanding the masses of colliding objects such as black holes or neutron stars is crucial. Two common parameterizations are used based on the analysis phase. The first combines the chirp mass $\mathcal{M}$, defined by     \[ \mathcal{M} = \left( \frac{(m_1 m_2)^{3/5}}{(m_1 + m_2)^{1/5}} \right), \cite{Cutler1994}\] ,

with the mass ratio $q = m_2 / m_1$ (assuming $m_2 \leq m_1$  and where $m_1$ and $m_2$ are the masses of the two objects), particularly useful during the inspiral phase for its influence on the binary orbit evolution and waveform. The second approach uses the individual masses $m_1$ and $m_2$ directly, more relevant in the ringdown phase where the total mass $M = m_1 + m_2$ predominates. This choice depends on how each parameter set affects the waveform's characteristics and the specific stage of the binary merger under consideration.


    \item \textbf{Spins}: The spin of each object in a binary system is an important factor for understanding the dynamics of the merger. Each object's spin is represented by a single three-dimensional vector that defines its angular momentum. This vector is characterized by both a magnitude, which represents the spin's strength, and by two angles (typically denoted as tilt and azimuthal angles) that define its orientation with respect to the orbital plane. Consequently, for a binary system comprising two objects, the total number of spin parameters is six: the magnitude, tilt angle, and azimuthal angle for each object's spin vector. These parameters affect the gravitational waveform's amplitude and phase, as well as the precession dynamics of the binary system.



    \item \textbf{Significance of Mass and Spin Estimates}: Estimating the masses and spins of objects such as black holes and neutron stars is fundamental for elucidating their formation mechanisms and evolutionary paths. Mass influences the frequency evolution of GW signals during the inspiral phase \cite{Blanchet2014}, and spin parameters affect both the amplitude modulation of the waveform throughout the coalescence \cite{Apostolatos1994} and the post-merger dynamics, including any potential recoil or 'kicks' experienced by the remnant black hole \cite{Campanelli2007, Campanelli2006b}. Such parameters are intrinsic to the GW source and provide insights into the energetic interactions at play.


   \item \textbf{Distance and Location}: The luminosity distance to a gravitational wave source informs us about the event's energy output, which help us to assess the scale of the phenomenon. The sky position allow for accurately locating the event, enhancing the potential for multi-messenger observations. These measurements are especially valuable in multi-messenger astronomy, allowing correlations between gravitational and electromagnetic signals to provide a more complete understanding of astrophysical events \cite{Multi-messenger Observations,Fairhurst2011}. Accurate distance and location data are essential for identifying host galaxies and contribute to our knowledge of the universe's expansion. Notably, while binary black hole (BBH) mergers typically do not yield electromagnetic counterparts, mergers involving neutron stars—such as binary neutron star (BNS) and neutron star-black hole (NSBH) mergers—do, due to their baryonic matter content \cite{Multi-messenger Observations}.


    \item \textbf{Orbital Parameters}: Orbital parameters, including inclination and eccentricity, inform our understanding of a binary system's dynamics as it progresses toward a merger. The orbital period, reflecting the time for one complete revolution, is not static but decreases as the orbiting objects emit gravitational waves and lose energy. Eccentricity, indicative of how much an orbit deviates from being circular, typically reduces over time, affecting the gravitational waveform's characteristics. Inclination can also vary, especially if the system exhibits precession — a phenomenon occurring when the spins are misaligned with the orbital angular momentum. Therefore, it is often more informative to discuss these parameters relative to a specific reference frequency, which provides a more consistent measure for gravitational wave analysis.


\end{itemize}


\subsubsection{Equation of State (EoS) and Its Importance}

One parameter of significant interest derived from GW signals, especially those from binary neutron star (BNS) mergers, pertains to the equation of state (EoS) of neutron stars \cite{EOS}. The EoS provides insights into the behavior of matter within the high-density environment of neutron stars. Understanding the EoS let us explore the characteristics of matter in conditions not reproducible in terrestrial laboratories, thereby advancing our knowledge of fundamental physics.


The EoS affects the GW signal in various ways:
\begin{itemize}

    \item \textbf{Tidal Deformability}: The EoS influences how much a neutron star is deformed by the gravitational field of its companion. This deformation leaves an imprint on the GW signal. Tidal deformability, denoted by \( \Lambda \), is a parameter that encapsulates this effect and is related to the EoS. It is defined as:
    \[ \Lambda = \frac{2}{3} k_2 \left( \frac{c^2 R}{Gm} \right)^5  \cite{Flanagan2008, Hinderer2008}\]
    Here, \( k_2 \) is the second Love number, which quantifies the star's response to the tidal potential of its companion \cite{Thorne1967, Binnington2009}. It is a dimensionless parameter that represents the ratio of the induced quadrupole moment of the neutron star to the external tidal field. The Love number depends on the internal structure of the star, which is governed by its EoS. Thus, measuring \( \Lambda \) from GW signals provides direct insight into the internal properties of neutron stars.


   \item \textbf{Post-Merger Signal}: The EoS also affects the post-merger signal in neutron star mergers. Different EoSs predict different outcomes, such as whether the remnant will immediately collapse into a black hole or form a hypermassive neutron star \cite{The-lifetime-of-binary-neutron-star-merger-remnants, EOS}.
\end{itemize}
 \textbf{These parameters shed light} on the ultra-dense matter's behavior, providing insights into nuclear physics and the state of matter in this scenarios.




\subsection{Bayesian Inference Framework in Gravitational-Wave Astronomy}
Bayesian inference is an important component of GW astronomy for analyzing and interpreting observational data. By employing Bayes' theorem, this framework facilitates the estimation of source properties, proving to be a substantial asset for PE. Tools such as the Bilby library exemplify the practical application of Bayesian inference in GW astronomy. Developed by Ashton, Hübner, Lasky, et al. \cite{Bilby: A user-friendly}, Bilby enhances the accessibility of Bayesian methods. Moreover, studies like those of Strub et al. \cite{Bayesian parameter-estimation of Galactic binaries}, which apply Bayesian inference to PE in LISA data for Galactic binaries, highlight the method's utility. Further advancements have been made by incorporating Bayesian inference with Markov Chain Monte Carlo (MCMC) techniques, leading to developments such as the Bilby-MCMC sampler by Ashton and Talbot \cite{Bilby-MCMC}.

\subsubsection{Fundamental Equation}
Bayesian inference in gravitational wave (GW) astronomy leverages Bayes' theorem \cite{Bilby: A user-friendly, Finn1992} to deduce the properties of GW sources from the data collected. This theorem is expressed as:
\begin{equation}
    p(\theta | d) = \frac{\pi(\theta)p(d | \theta)}{p(d)}
\end{equation}
where:
\begin{itemize}
    \item $p(\theta | d)$ denotes the posterior probability density function (PDF) for the parameters $\theta$ given the data $d$.
    \item $\pi(\theta)$ is the prior probability distribution of the parameters.
    \item $p(d | \theta)$ represents the likelihood function, which is the probability of observing the data $d$ given the parameters $\theta$.
    \item $p(d)$, known as the Bayesian evidence, acts as the normalization constant.
\end{itemize}

\subsubsection{Likelihood Function}
For GW data with the assumption of Gaussian noise, the likelihood function is formulated as:
\begin{equation}
    \ln L(d | \theta) = -\frac{1}{2} \sum_{k} \left( \frac{[d_k - h_k(\theta)]^2}{\sigma_k^2} + \ln 2\pi\sigma_k^2 \right)
\end{equation}
In this equation:
\begin{itemize}
    \item $d_k$ indicates the observed data at frequency bin $k$.
    \item $h_k(\theta)$ is the expected signal model as a function of the source parameters $\theta$.
    \item $\sigma_k$ describes the noise standard deviation in each frequency bin.\cite{Smith2016, Wong2023}.
\end{itemize}

\
\subsection{Fisher Information Matrix in GW Parameter Estimation}
The Fisher Information Matrix (FIM) serves as an analytical instrument in GW astronomy for PE. Within the framework of frequentist statistics, the FIM measures how much information an observable variable provides about the parameters of a model, which in turn informs the precision limits of these measurements. Thorne's work "Bayesian inference for gravitational waves and pulsars" \cite{Thorne2002} comprehensively reviews GW parameter estimation, with an emphasis on the FIM's utility in both frequentist and Bayesian methodologies. The FIM is particularly useful in scenarios with high signal-to-noise ratios (SNR), contributing to the accurate estimation of source characteristics, including the masses and spins of compact objects within binary systems.

\subsubsection*{Definition and Importance}
The Fisher Information Matrix (FIM) elements are determined through the inner product:
\begin{equation}
    \Gamma_{\mu\nu} = \left( \frac{\partial h}{\partial \theta_\mu}, \frac{\partial h}{\partial \theta_\nu} \right)
\end{equation}
where $\theta_\mu$ and $\theta_\nu$ denote distinct parameters of the gravitational wave (GW) signal. The inner product is described as:
\begin{equation}
    (g, h) = 4 \text{Re} \int_{0}^{\infty} \frac{\tilde{g}(f) \tilde{h}^*(f)}{S_n(f)} df \cite{Cutler1994}
\end{equation}
In this formulation:
\begin{itemize}
    \item $\tilde{g}(f)$ and $\tilde{h}(f)$ represent the Fourier transforms of the time-domain functions $g(t)$ and $h(t)$, respectively.
    \item $S_n(f)$ is the power spectral density of the noise in the detector.
\end{itemize}


\textbf{The Fisher Information Matrix (FIM)} effectiveness compared to real parameter estimation techniques, such as Markov-chain Monte Carlo (MCMC), has been examined in the context of gravitational-wave signals from merging compact binaries. Rodriguez, Farr, Farr, and Mandel \cite{Rodriguez2013} found that the FIM can overestimate the uncertainty in PE achievable by MCMC, particularly for systems composed of binary black holes. Their study highlights the increasing discrepancies in the uncertainty estimates provided by FIM and MCMC as the total mass of the binary system increases.


\subsection{Markov Chain Monte Carlo Methods}
MCMC methods, and specifically the Metropolis-Hastings (MH) algorithm, are utilized for sampling from the posterior distribution, effectively navigating the parameter space to estimate gravitational wave sources' properties \cite{vanderSluys2008}. The MH algorithm operates through an iterative process, proposing new parameter states and determining their acceptance based on predefined criteria:

\begin{enumerate}
    \item Initialize with a starting point \( \theta^{(0)} \).
    \item For each iteration \( i \), propose a new parameter state \( \theta' \) from a proposal distribution \( q(\theta' | \theta^{(i-1)}) \).
    \item Calculate the acceptance probability \( \alpha \) as:
    \[ \alpha(\theta', \theta^{(i-1)}) = \min\left(1, \frac{p(\theta' | D) q(\theta^{(i-1)} | \theta')}{p(\theta^{(i-1)} | D) q(\theta' | \theta^{(i-1)})}\right) \]
    \item Accept \( \theta' \) with probability \( \alpha \), otherwise retain \( \theta^{(i-1)} \).
    \item Repeat for a significant number of iterations.
\end{enumerate}

In this context, \( q(\theta' | \theta^{(i-1)}) \) represents the proposal distribution, which is a mechanism for generating new candidate states \( \theta' \) based on the current state \( \theta^{(i-1)} \). This distribution is crucial for the exploration of the parameter space, guiding the sampling process towards areas of higher posterior probability.

\subsection{Computational Challenges in GW Parameter Estimation}


\begin{itemize}


\item \textbf{Complexity of High-Dimensional Space}: GW signals depend on numerous parameters that define a multi-dimensional parameter space. As detailed in Table 1, a typical binary system's parameter space might extend beyond 15 dimensions, with the vectorial nature of spins contributing to the dimensional complexity. Efficiently exploring this space for parameter estimation (PE) requires robust computational techniques. Reduced-order quadrature and adaptive proposal distributions are effective in managing these dimensions \cite{Smith2016}. Additionally, the use of hardware accelerators like GPUs enhances the computation of waveform models and likelihood evaluations, facilitating faster processing in extensive parameter landscapes \cite{Wong2023}.






\item \textbf{Sampler Efficiency}: In GW astronomy, the efficiency of sampling methods, notably Markov Chain Monte Carlo (MCMC), is essential for navigating complex parameter spaces \cite{LISA-MCMC}. Sampler efficiency relates specifically to the exploration of these spaces to derive parameter estimates. MCMC methods face challenges such as:
    \begin{itemize}
        \item \textit{Multimodal Distributions}: In GW parameter spaces, the probability density functions may have multiple peaks corresponding to viable parameter sets. Efficient sampling across these modes is necessary, especially when probability valleys separate them \cite{Roulet2022}.
        \item \textit{Correlations Between Parameters}: Parameters such as the masses and spins of black holes might be correlated, influencing the observed gravitational waveforms. These correlations can lead to degenerate regions in parameter space, complicating the sampling process and potentially delaying convergence \cite{Luengo2020, Kramer2014}.
    \end{itemize}




These challenges extend beyond the mere existence of numerous dimensions; they specifically relate to the complexity of the parameter space that must be navigated during sampling. Overcoming the limitations of traditional MCMC methods in gravitational wave parameter estimation (GW PE) necessitates the adoption of advanced sampling strategies. Techniques such as heterodyned likelihood and normalizing flows, discussed by Wong et al. \cite{Wong2023}, improve sampling efficiency by streamlining likelihood evaluation and navigating complex parameter spaces more effectively. Moreover, gradient-based samplers exploit derivative information of the posterior to steer the sampling towards regions of higher probability, enhancing efficiency in exploring multidimensional spaces \cite{Wong2023}.
Alternative sampling methods are crucial for managing complex distributions and resolving parameter correlations to ensure precise and reliable parameter estimation.
\item \textbf{Simulated Annealing}: This method modifies the likelihood surface to facilitate exploration, potentially accelerating convergence. Critical to its success is the correct selection of initial temperature and annealing parameters, which requires meticulous tuning to optimize performance \cite{Wong2023}. Simulated Annealing's flexibility makes it a viable option for addressing the intricate parameter spaces typical in GW data analysis.



    \item \textbf{Future Directions}: The increase in GW observations and the expanding volume of data necessitate ongoing improvements in computational methods. Investigating the application of machine learning for data analysis, refining waveform models for enhanced precision, and utilizing progressive computational software are among the principal areas for future exploration and development in this domain \cite{Wong2023}.


\end{itemize}

\clearpage

    \begin{figure}[ht]
    \centering
    \includegraphics[width=0.99\textwidth]{ParaEstimExample.png}
    \caption{The detector strain measurement of GW150914 (B.abbot et al., 2016c) \cite{B.abbot et al2016c} observed by the LIGO Hanford(H1,Top) and Livingston(L1, bottom) detectors in gray. Times are realative to 14/9/2015 at 09:50:45 UTC. The cyan curve is the estimated signal using IMRPhenom and EOBNR waveform template, the blue curve a signal reconstruction based on BayseWave.\cite{B.abbot et al2016d}}
    \label{fig:gw_recovery}
    \end{figure}




Figure \ref{fig:gw_recovery} illustrates the comparison between the detector strain measurement of GW150914 by LIGO Hanford and Livingston and the modeled GW signals. The graph displays the complexities involved in accurately matching the observed detector data, represented in gray, with theoretical models. The discrepancies between these curves highlight the intricate task of waveform reconstruction and PE in gravitational wave astronomy.


\section{Research Objective}
The detection of a gravitational wave (GW) event initiates parameter estimation (PE), which underpins the astrophysical interpretation of the source and any subsequent test of general relativity. The LIGO--Virgo--KAGRA collaborations deliver PE results in two distinct stages. A low-latency PE run, executed within tens of minutes of the trigger on the preliminary noise PSD \cite{Chaudhary2024, Bilby: A user-friendly}, drives multi-messenger alerts and astronomer follow-up. A second, calibration-finalized PE run is performed weeks to months later for inclusion in catalog releases and population analyses \cite{GWTC3}. The shift in the PSD between these two stages is small: Sun et al.\ report systematic calibration errors below $\sim 2\%$ in amplitude and $\sim 2^\circ$ in phase across the most sensitive band \cite{Sun2020}, and Biscoveanu et al.\ show that propagating the resulting PSD uncertainty into the posteriors changes the inferred parameters by only a few percent \cite{Biscoveanu2020}. Yet the final PE run is launched from scratch, re-sampling the full 15-dimensional astrophysical prior despite the low-latency posterior having already constrained the bulk of the high-probability region. This duplication is the bottleneck this thesis addresses.

The primary goal of this research is to accelerate the final, catalog-quality PE run by re-using the information already contained in the low-latency posterior. We propose a two-phase method. Phase 1 (Scout) is functionally analogous to the existing low-latency PE — a nested sampler on the preliminary PSD with uninformative priors — and yields an approximate posterior. Phase 2 (Refine) launches the final, calibration-finalized PE not from the broad astrophysical prior but from a truncated Gaussian prior constructed from the Phase 1 posterior. By restricting the sampler to the already-identified high-probability region while retaining the same final-PSD likelihood evaluation, Phase 2 converges substantially faster than a blind run. The hypothesis under test is that the resulting Bayes factor and posterior agree with the blind baseline within statistical noise, at a fraction of the wall-clock cost.

\subsection{ Research expectation }

Based on our preliminary analysis with simplified 6-parameter runs (Table~\ref{tab:compute_specs}), we project significant computational savings for the final, calibration-quality 15-parameter PE run. Although our test runs with 6 parameters complete in $\sim$3 hours, production-grade final PE with 15 parameters typically requires 20–30 hours per event \cite{Romero-Shaw2020}. Our two-phase approach replaces the blind final run with a prior-compressed Refine pass that re-uses the low-latency posterior; we estimate this Refine pass will require approximately 56\% (see Appendix~\ref{appendix:Computational Time Savings Analysis}) of the blind run's wall-clock cost. The savings derive from
\begin{itemize}
    \item a reduced effective prior volume after truncated-Gaussian compression,
    \item accelerated convergence from informed proposal distributions, and
    \item reduced sampler ``burn-in'' relative to an uninformative-prior start.
\end{itemize}

For a typical observing run with $\sim$200 events, this could reduce the CPU cost of producing catalog-quality posteriors from approximately 4{,}200 hours (175 days) to 2{,}100 hours (87.5 days), a saving that compounds across waveform families and population re-analyses. This 44\% reduction enables more thorough analyses, additional waveform models, and faster catalog data releases.

    \begin{figure}[ht]
    \centering
    \includegraphics[width=0.50\textwidth]{noise_curve_importance.png}
\caption{The plot illustrates the strain sensitivity of the Advanced LIGO detectors during the initial (O1) and the final science run (S6) of the initial LIGO detectors. The dark red and light red curves represent the O1 strain noise for H1 and L1 respectively, indicating comparable performance between the two detectors. Also depicted are the Advanced LIGO design sensitivity and a prospective future upgrade \cite{Miller2015, Abbott2016-sensitivity}, underscoring the potential for discovery in the upcoming years.}

    \label{fig:noise_curve_importance}
    \end{figure}
Figure~\ref{fig:noise_curve_importance} contextualizes the role of detector sensitivity in PE accuracy. Our research tests whether the posterior obtained from a Scout run on a preliminary PSD is sufficiently close to the truth that it can seed an accelerated final PE run on the calibration-finalized PSD. The figure illustrates that stable noise curves improve detection precision, motivating the question of how much information from the preliminary-PSD Scout run is preserved when the PSD is updated for the final analysis.

%\newpage

\section{Research Assumptions and/or hypotheses }

\begin{itemize}
 \item \textbf{Re-use of low-latency posteriors:} The current pipeline performs two PE runs per event — a low-latency run on the preliminary PSD and a final run on the calibration-finalized PSD — but discards the low-latency posterior when launching the final run. This research tests whether constructing the final run's prior from the low-latency (Scout) posterior accelerates convergence without biasing the final result. The method assumes that the Scout posterior contains enough information about the high-probability region of parameter space that the Refine sampler need only explore a compressed prior around it, rather than the full astrophysical prior.
  \begin{itemize}
    \item \emph{Validation Method:} Time saved and the precision of PE will be quantitatively analyzed.
  \end{itemize}



    \item \textbf{Hypothesis on Gradient Descent Versus MCMC:} The exploration of gradient descent optimization alongside Metropolis-Hastings MCMC methods considers the contrasting demands of computational efficiency and the ability to navigate complex parameter landscapes. Gradient descent, often more straightforward and less computationally intensive, is advantageous in scenarios where the posterior distribution is smooth and unimodal, allowing for rapid convergence to the maximum likelihood. However, in a multi-dimensional parameter space common in gravitational wave astronomy, this simplicity may be offset by MCMC's robustness in exploring complex, potentially multimodal distributions where gradient descent might falter.

  \begin{itemize}
    \item \emph{Empirical Testing:} The hypothesis will be empirically tested by comparing the convergence times and the fidelity of the parameter estimate in a diverse set of GW data.
    \item \emph{Scope note (added at submission):} A direct gradient-descent vs.\ nested-sampling benchmark was de-scoped from the present work in favour of completing the systematic SNR sweep and the matched-budget control experiment. The prior-compression motivation for the gradient-descent analogy is discussed qualitatively in Section~\ref{sec:theoretical_interpretation}; the direct empirical comparison is identified as a concrete follow-up in Section~\ref{sec:future_work}.
  \end{itemize}
\end{itemize}

These hypotheses will be subjected to rigorous empirical testing, with outcomes informing the preferred methodological approach for real-time GW analysis in future observational runs.



%\newpage

\section{Research Plan}
\subsection*{Overview and Theoretical Framework}

This research focuses on optimizing parameter estimation (PE) processes in gravitational wave (GW) astronomy using Bayesian inference and MCMC methods within the Bilby framework \cite{Bilby-MCMC, Bilby: A user-friendly}. We aim to enhance the efficiency and accuracy of PE by employing adaptive proposal distributions and convergence diagnostics. Our strategy includes a two-step approach of using preliminary estimates from early noisy data and refining these with cleaner data, aiming to reduce the timeline for PE and expedite scientific outputs.




\subsection{Optimization of MCMC Sampling}

This research focuses on refining the efficiency of Markov Chain Monte Carlo (MCMC) sampling for gravitational wave parameter estimation. Central to this refinement is the optimization of the proposal distribution, which dictates the selection probabilities of parameter values during sampling. We plan to implement advanced tuning techniques that adapt the proposal distribution based on performance metrics gathered during sampling. This includes adjusting the step sizes or adopting more sophisticated algorithms like adaptive Metropolis or Hamiltonian Monte Carlo if the parameter space shows complex dependencies or multimodality. Additionally, convergence diagnostics will be systematically employed to verify the thorough exploration of the parameter space and validate the sampling efficiency. These diagnostics will involve statistical tests such as the Gelman-Rubin statistic and effective sample size calculations to confirm that the chains are mixing well and converging to the target distribution. Supported by the capabilities of the Bilby framework, these strategies are instrumental in enhancing the accuracy and reliability of our parameter estimation models \cite{Bilby-MCMC, Smith2016}.

\subsubsection{Advanced Tuning Techniques and Convergence Diagnostics}

In the optimization of MCMC sampling, advanced tuning techniques will be crucial \cite{Luengo2020}. These involve algorithmically adjusting the step sizes and shapes of the proposal distribution based on the performance metrics of the MCMC samples. This process, termed `adaptive MCMC',\cite{Haario2001, Roberts2009} aims to enhance the exploration of the parameter space by dynamically updating the proposal distribution to reduce correlation between consecutive samples and avoid local optima \cite{Kramer2014}.

Additionally, we will utilize rigorous convergence diagnostics to verify the effectiveness of the sampling strategy. These diagnostics include the calculation of the Gelman-Rubin statistic \cite{GelmanRubin1992}, which compares the variance between multiple MCMC chains to the variance within each chain to assess convergence. Another important diagnostic tool is the Effective Sample Size (ESS), which measures the number of independent samples in the chain, providing insights into the chain's efficiency in exploring the parameter space \cite{vanderSluys2008}. Both tools are essential for confirming that the chains have sufficiently explored the parameter space and are representative of the true posterior distribution.

These advanced techniques and diagnostics are integral to our methodology and are fully supported by the Bilby framework, which provides built-in functions for implementing these adaptive strategies and evaluating their performance \cite{Bilby-MCMC, Smith2016}.

\subsubsection{Adaptive Proposal Distribution}
The modification of proposal distributions during MCMC sampling affects the sampling efficiency. Our method implements algorithms that update the proposal distribution based on the MCMC chain data. This adaptation aims to reduce convergence time and improve parameter space exploration through continuous updates of the proposal distribution during sampling. These methods have improved MCMC performance in tests with the Bilby framework \cite{Bilby-MCMC, Bilby: A user-friendly}.

\subsubsection{Convergence Diagnostics}

Reliable convergence diagnostics are fundamental for validating the MCMC sampling process. We plan to apply the autocorrelation function (ACF) to measure the independence of the samples within the chains. The ACF at lag \( k \) helps identify the point at which samples become uncorrelated, calculated by:
\[
ACF(k) = \frac{\sum_{t=1}^{N-k} (X_t - \overline{X})(X_{t+k} - \overline{X})}{\sum_{t=1}^N (X_t - \overline{X})^2},\cite{Geyer1992, Sokal1997},
\]
with \( X_t \) representing the samples and \( \overline{X} \) their mean. Additionally, the potential scale reduction factor (PSRF) will be employed to assess chain mixing and convergence, ensuring that multiple chains are consistent with one another:
\[
PSRF = \frac{n - 1}{n} + \frac{m+1}{mn} \frac{B}{W}, \cite{GelmanRubin1992}, \]
where \( B \) is the between-chain variance, \( W \) is the within-chain variance, \( m \) is the number of chains, and \( n \) is the chain length. These techniques, crucial for monitoring the adequacy of our sampling process, align with established methods within the Bilby framework \cite{Bilby-MCMC}.


\subsection{Parameter Estimation:Testing Gradient Descent Optimization}

On the other end, the task of PE can be approached using gradient descent optimization. While MCMC with the Metropolis-Hastings algorithm provides a robust mechanism for sampling from the posterior distribution, gradient descent offers a simpler and potentially faster path to optimization, particularly when the likelihood function is differentiable with respect to the parameters.

\begin{enumerate}
    \item Begin with initial parameter estimates derived from the degraded noise signal, a necessary step to establish a starting point for optimization.
    \item Apply gradient descent optimization to iteratively refine the parameter estimates. This method adjusts parameters by stepping in the direction of the steepest descent of the likelihood function, mathematically described as:
    \[
    \theta^{(new)} = \theta^{(old)} - \eta \nabla_\theta \mathcal{L}(\theta^{(old)}) , \cite{Boyd2004, Nocedal2006}
    \]
    where \( \theta^{(old)} \) and \( \theta^{(new)} \) are the old and new parameter vectors, respectively, \( \eta \) is the learning rate, and \( \nabla_\theta \mathcal{L} \) is the gradient of the likelihood function \( \mathcal{L} \) with respect to \( \theta \).

\end{enumerate}

\subsection{Simulation and Validation}

The simulation and validation phase tests the signal processing and PE algorithms. We test if the algorithms work correctly with actual detector data by following these steps:

\begin{enumerate}
    \item \textbf{Development of Simulated Signals:} We generate a suite of synthetic GW signals and noise that closely mimic the data obtained by detectors like LIGO and Virgo. This simulation involves:
    \begin{itemize}
        \item Creating realistic GW waveforms based on astrophysical models of binary mergers.
        \item Superimposing these signals onto synthetic noise backgrounds that reflect the statistical properties of actual detector noise.
    \end{itemize}

    \item \textbf{Signal Degradation and Noise PSD Adjustment:}
    To evaluate our two-phase method's potential for rapid and efficient PE, we simulated varying noise conditions. Rather than strictly degrading the injected signal itself, we adjust the noise Power Spectral Density (PSD) used during parameter estimation. This method allows the PE process (conducted using Bilby) to account for different noise conditions by specifying various noise PSD files.

    The process of noise refinement is applied by employing a moving average filter to the frequency domain of the noise data, following standard signal processing techniques \cite{Welch1967}. This technique adjusts the noise data to represent varying observational conditions, enabling a more stable and consistent noise curve. The refinement is controlled by a noise time factor, where an increased factor leads to a smoother, albeit less detailed, noise curve, simulating the different noise levels encountered during actual observations (e.g., early-run preliminary calibration versus late-run finalized calibration).

    Figure \ref{fig:noise_refinement_comparison} demonstrates the varying Amplitude Spectral Densities (ASDs) generated for the LIGO H1 detector across different temporal windows (from 20 minutes to 1152 hours). These variations serve as the testing ground for applying initial calibrations (Phase 1) and subsequent iterative refinements (Phase 2).

    \begin{figure}[htbp]
        \centering
        \begin{subfigure}[b]{0.45\textwidth}
            \includegraphics[width=\textwidth]{H1_noise_ASD_0.333_hours.png}
            \caption{ASD for 20 minute window}
            \label{fig:H1_250}
        \end{subfigure}
        \hfill
        \begin{subfigure}[b]{0.45\textwidth}
            \includegraphics[width=\textwidth]{H1_noise_ASD_24.167_hours.png}
            \caption{ASD for 24 hours window}
            \label{fig:H1_500}
        \end{subfigure}

        \vspace{0.5cm}
        \begin{subfigure}[b]{0.45\textwidth}
            \includegraphics[width=\textwidth]{H1_noise_ASD_168.167_hours.png}
            \caption{ASD for 168 hours window}
            \label{fig:H1_750}
        \end{subfigure}
        \hfill
        \begin{subfigure}[b]{0.45\textwidth}
            \includegraphics[width=\textwidth]{H1_noise_ASD_1152.166_hours.png}
            \caption{ASD for 1152 hours window}
            \label{fig:H1_1000}
        \end{subfigure}

        \caption{Comparison of ASD at different noise windows on the LIGO H1 detector. Smoother curves represent longer integration times, simulating finalized noise calibration.}
        \label{fig:noise_refinement_comparison}
    \end{figure}

    \begin{figure}[ht]
        \centering
        \includegraphics[width=0.45\textwidth]{noise_example.jpg}
        \caption{An example of simulated signal (green) with added noise (orange) to simulate an early detection strain (blue) as opposed to the more certain strain which detected at O3 which was taken from the Bilby documentation\cite{Bilby: A user-friendly}}
        \label{fig:noise_example}
    \end{figure}

    \item \textbf{Running the Algorithm}
    \label{sec:running_algorithm}
    This stage involves applying our algorithms to 'start-of-run' data, which is characterized by higher noise levels and serves as our 'degraded' dataset. The objective is to extract preliminary parameter estimates from this challenging data. We will then refine these estimates using the 'end-of-run' data, which corresponds to the actual, cleaner signals. This process aims to replicate a real-world scenario where early-run noise profiles are used for initial estimation, and end-of-run data, with improved noise characteristics, are used for final PE.

    \begin{enumerate}
      \item Processing 'start-of-run' noisy data to obtain preliminary parameter estimates.
      \item Refining estimates with 'end-of-run' data, reflecting cleaner signal conditions.
      \item Evaluating the computational efficiency and accuracy of this two-step estimation process.
    \end{enumerate}

    The success of this two-stage estimation process will be gauged by the accuracy of parameters obtained and the reduction in computational time and resources.

    \item \textbf{Algorithm Refinement}
    Algorithm refinement includes adjusting algorithm parameters, such as proposal distribution variances (via Truncated Gaussians), to optimize performance. Validating the refined parameter estimates against the known parameters of the simulated signals ensures that the optimization is accurately capturing the true parameter values without introducing statistical bias.

\end{enumerate}

\subsection{Validation and Implementation}
This research concludes with testing of the MCMC and gradient optimization algorithms. We will assess time reduction using past gravitational wave events, focusing on the measurement of main parameters including masses, spins, and sky locations. The tests will examine how the algorithms reduce delays in the estimation pipeline, which limits the speed of gravitational wave data analysis. We will analyze the computational costs to verify that time savings do not require excess resources. After validation, the algorithms can be added to the analysis process for future gravitational wave observations, improving analysis speed. This improvement will help connect data collection to analysis results more quickly, supporting the combination of gravitational wave and electromagnetic observations.
To validate our 44\% refinement time estimate, we will:
\begin{itemize}
    \item Compare full PE versus two-phase PE on a subset of real GW events from O3
    \item Measure actual refinement convergence rates with different Phase 1 starting points
    \item Test sensitivity to noise curve evolution patterns
    \item Document the relationship between initial estimate quality and refinement time
\end{itemize}

\subsubsection*{Distributed Computing for Large-scale Inference}

Given the computationally intensive nature of MCMC methods, the research will leverage distributed computing resources, utilizing high-performance computing clusters and cloud platforms for parallel processing of MCMC chains.
For this work I have used my personal PC-  full spec in the appendix A.3

\subsubsection{Validation of Parameter Recovery Accuracy}

A critical component of our validation strategy involves quantifying the accuracy difference between the two-phase approach and traditional single-phase PE, particularly for astrophysically significant parameters. We will implement systematic comparisons focusing on:

\begin{itemize}
    \item \textbf{Sky localization accuracy:} The sky position parameters (RA, Dec) directly impact electromagnetic follow-up observations. We will quantify the sky area uncertainty (typically measured in square degrees at 90\% confidence) for both methods and assess whether the two-phase approach maintains sufficient accuracy for successful multi-messenger astronomy.

    \item \textbf{Mass ratio determination:} The mass ratio ($q$) provides crucial insights into binary formation channels and population synthesis models. We will evaluate the systematic bias and statistical uncertainty in mass ratio recovery, ensuring that the two-phase method preserves the precision needed for astrophysical inference.

    \item \textbf{Quantitative metrics:} For each parameter, we will compute:
    \begin{equation}
    \Delta_{\text{param}} = \frac{|\theta_{\text{two-phase}} - \theta_{\text{single-phase}}|}{\sigma_{\text{single-phase}}}
    \end{equation}
    where $\theta$ represents the median posterior value and $\sigma$ the 1-$\sigma$ uncertainty. Success criteria require $\Delta_{\text{param}} < 1.0$ for all critical parameters. The initial target of $\Delta_{\text{param}} < 0.5$ was set in anticipation of high-SNR operation only; once the systematic evaluation was extended to moderate and marginal detections (150--300\,Mpc, $\rho \lesssim 30$), the posterior widths broaden and the $1.0$ criterion --- Phase 2 median falling within the Baseline $1\sigma$ credible interval --- is the appropriate statistical standard for the full SNR range.
\end{itemize}

This validation will be performed on a representative sample of O3 events spanning different signal-to-noise ratios and binary configurations to ensure robust performance across the full parameter space encountered in gravitational wave astronomy.

\subsubsection{Parallelization Strategy and Wall-Clock Time Analysis}

While our analysis projects 2,100 CPU hours for Phase 2 refinement across ~200 events, the actual wall-clock time depends critically on parallelization efficiency. Modern computational infrastructure enables several levels of parallelization:

\textbf{Event-level parallelization:} Each GW event can be analyzed independently, allowing perfect linear scaling up to 200 parallel processes. With access to a computing cluster with 200 cores, the entire catalog could theoretically be processed in:
\begin{equation}
t_{\text{wall}} = \frac{2,100 \text{ CPU hours}}{200 \text{ cores}} = 10.5 \text{ hours}
\end{equation}

\textbf{Chain-level parallelization:} Within each event, MCMC chains can run in parallel. Typical configurations use 4-8 parallel chains per event, reducing single-event processing from ~10.5 hours to 1.5-2.5 hours with sufficient resources.

\textbf{Practical considerations:} Real-world constraints include:
\begin{itemize}
    \item Queue waiting times on shared computing resources
    \item Memory bandwidth limitations when running many parallel jobs
    \item I/O bottlenecks when multiple processes access data simultaneously
    \item Load balancing inefficiencies (some events converge faster than others)
\end{itemize}

Based on typical LIGO/Virgo computing infrastructure with 1000 available cores and accounting for practical overhead, we estimate that the Phase 2 refinement for 200 events could be completed in 24--48 hours of wall-clock time. This represents a substantial improvement over the conventional approach, in which the final, calibration-quality run is launched from the broad astrophysical prior and discards the low-latency posterior already obtained at trigger time.

Future work will include benchmarking on actual computing clusters to refine these estimates and optimize the parallelization strategy for different computational architectures. We will also investigate GPU acceleration for likelihood evaluations, which could further reduce both CPU hours and wall-clock time by factors of 10-100 for suitable portions of the calculation \cite{Bilby: A user-friendly}.

\section{Preliminary Results}

\subsection{Noise Refinement Process}
In this research, the process of noise refinement is applied by employing a moving average filter to the frequency domain of the noise data, following standard signal processing techniques \cite{Welch1967}. This technique adjusts the noise data to represent varying observational conditions, enabling a more stable and consistent noise curve. Such refinement is crucial for improving the detectability of gravitational wave signals and the accuracy of the parameter estimation. The refinement is controlled by the noise time factor, where an increased factor leads to a smoother, albeit less detailed, noise curve, simulating different noise levels that might be encountered during actual observations.

\subsection{Analyzing the Influence of Noise on Parameter Estimation}
The preliminary results presented here highlight the impact of different noise levels on the accuracy of PE within gravitational-wave data analysis. This section demonstrates our method's potential for rapid and efficient PE through the application of initial calibrations and subsequent iterative refinements, aligning with our goal of efficient computational resource management.

\begin{figure}[htbp]
    \centering
    \begin{subfigure}[b]{0.45\textwidth}
        \includegraphics[width=\textwidth]{H1_noise_ASD_0.333_hours.png}
        \caption{ASD for 20 minute window}
        \label{fig:H1_250}
    \end{subfigure}
        \begin{subfigure}[b]{0.45\textwidth}
        \includegraphics[width=\textwidth]{H1_noise_ASD_24.167_hours.png}
        \caption{ASD for 24 hours window}
        \label{fig:H1_500}
    \end{subfigure}
        \begin{subfigure}[b]{0.45\textwidth}
        \includegraphics[width=\textwidth]{H1_noise_ASD_168.167_hours.png}
        \caption{ASD for 168 hours window}
        \label{fig:H1_750}
    \end{subfigure}
        \begin{subfigure}[b]{0.45\textwidth}
        \includegraphics[width=\textwidth]{H1_noise_ASD_1152.166_hours.png}
        \caption{ASD for 1152 hours window}
        \label{fig:H1_1000}
    \end{subfigure}


    \caption{Comparison of ASD at different noise windows on the LIGO H1 detector at different integration times.
    }
    \label{fig:noise_refinement_comparison}
\end{figure}

\subsection{Technical Discussion}
The technical efficacy of our approach can be observed in the noise spectrum comparisons and PE plots. These visual representations underscore how noise calibration at different intervals significantly alters the outcome of the PE. Our methodology's strength lies in the utilization of intermediate results obtained from preliminary noise analysis to inform the final estimation process. By implementing an adaptive strategy that capitalizes on the incremental nature of noise refinement and PE, we aim to establish a more streamlined and efficient pipeline for analyzing gravitational wave data.

Through this ongoing research, we will further explore the quantitative differences between PE outcomes under varying noise conditions and develop an algorithm that employs initial estimates to curtail the runtime for subsequent analyses with final noise calibration. The anticipated outcome is a substantial reduction in the time required to achieve final parameter estimates, facilitating a faster turnaround for scientific interpretation and communication of findings within the scientific community.

\subsection{Analysis Duration and Scaling Considerations}
Parameter estimation analyzes were performed on the system described in Table~\ref{tab:compute_specs}. The total computation time for the six-parameter estimate was 2.66 hours, with an average chain length of 10,517 samples per parameter. This computational cost highlights the need for our proposed methods that aim to reduce the time needed for complete PE after the final noise calibration becomes available.

Table ~\ref{tab:compute_specs} shows the runtimes for simplified 6-parameter PE. Production PE includes:
 \begin{itemize}
      \item 2 masses parameters (chirp mass and mass ratio).
      \item 6 spin parameters (3D vectors for each object)
      \item 4 extrinsic parameters (distance, sky location, inclination, phase).
    \item 3 additional parameters (time, polarization, etc.)

    \end{itemize}
Total: 15 parameters

The computational cost scaling from 6 to 15 parameters is approximately:
\begin{itemize}
\item MCMC sampling efficiency decreases with dimension
\item More iterations needed for convergence
\item Expected scaling factor: $\sim 7\times$ (derived below)
\end{itemize}

The scaling factor follows from the \texttt{dynesty} cost relation $N_{\rm eval} \propto n_{\rm live} \cdot n_{\rm walks} \cdot \ln(V_{\rm prior}/V_{\rm posterior})$, where the logarithmic compression factor grows approximately linearly with dimensionality \cite{Skilling2006,HigsonNS2019,Romero-Shaw2020}: $\ln \sim 4$ at $d=6$ rises to $\ln \sim 14$ at $d=15$, a ratio of $\sim 3.5\times$. An additional $\sim 2\times$ accrues from the higher per-likelihood cost of precessing waveforms (\texttt{IMRPhenomXPHM} vs.\ \texttt{IMRPhenomD}) \cite{Pratten2020}, yielding a combined factor of $\sim 7\times$ consistent with the empirical 6D-to-15D scaling reported by Romero-Shaw et al.\ \cite{Romero-Shaw2020}. Therefore, our $\sim$3 hour test runs project to 20--30 hour production runs, making the estimated 50 percent time savings even more significant at scale.


\subsection{Future Work and Directions}
The next stages of this research will broaden the range of parameters assessed for their impact on gravitational waveforms. Key to this expansion will be the quantitative evaluation of how initial noise conditions versus refined conditions affect the outcomes of parameter estimation. An algorithm will be developed to utilize the results obtained under initial noise conditions to reduce the computation time required for subsequent analyses with refined noise. This strategy aims to merge preliminary and refined estimates to streamline the parameter estimation process, enhancing the efficiency of our method.




%\clearpage


\newpage

\section{Results}
\label{sec:results}

\subsection{Baseline Configuration and Computational Scaling}
\label{sec:baseline}
To establish a baseline for computational efficiency, initial parameter estimation (PE) analyses were performed on a simplified 6-parameter space using the system configuration detailed in Appendix A. The total computation time for the six-parameter estimate was 2.66 hours, achieving an average chain length of 10,517 samples per parameter.

While a 6-parameter run provides fundamental convergence metrics, a production-level gravitational-wave parameter estimation requires mapping a 15-dimensional space. The complete parameter space includes:
\begin{itemize}
    \item 2 mass parameters (chirp mass and mass ratio).
    \item 6 spin parameters (3D vectors for each object).
    \item 4 extrinsic parameters (distance, sky location, inclination, phase).
    \item 3 additional parameters (geocentric time, polarization, etc.).
\end{itemize}

The computational cost of nested sampling scales with dimensionality through two channels: the per-likelihood waveform-evaluation cost (roughly independent of $d$) and the number of likelihood evaluations $N_{\rm eval}$ required for convergence. For \texttt{dynesty},
\begin{equation}
    N_{\rm eval} \;\propto\; n_{\rm live} \cdot n_{\rm walks} \cdot \ln\!\left(\frac{V_{\rm prior}}{V_{\rm posterior}}\right),
\end{equation}
where the logarithmic compression factor is itself approximately linear in the dimensionality $d$, since each additional well-constrained parameter contributes an additive $\ln(\sigma_{\rm prior}/\sigma_{\rm posterior})\sim 3$--$5$ to the sum \cite{Skilling2006,HigsonNS2019,Romero-Shaw2020}. For $d_1 = 6 \to d_2 = 15$, the prior-volume compression factor therefore grows from $\sim\!50$ ($\ln \sim 4$) to $\sim\!10^6$ ($\ln \sim 14$), a ratio of $14/4 \approx 3.5$. The additional $\sim 2\times$ comes from increased per-likelihood cost when the additional parameters include precessing spins (\texttt{IMRPhenomXPHM} vs.\ \texttt{IMRPhenomD} has $\sim 2$--$3\times$ the evaluation cost, \cite{Pratten2020}). Combining the two yields $N_{\rm eval}^{(15)} / N_{\rm eval}^{(6)} \approx 3.5 \times 2 \approx 7$, consistent with the empirical scaling reported in \cite{Romero-Shaw2020} (Table~3) for transitions between fixed-spin and full-spin configurations in similar nested-sampling runs. Extending the $\sim 2.7$-hour 6D test runs to the full 15-parameter space therefore projects to $\sim 19$--$27$ hours of continuous computation per event, with the upper end of the range assumed in the catalog-scale estimate of Section~\ref{sec:catalog}. This significant computational burden establishes the baseline necessity for the two-phase estimation method, aimed at reducing the overall wall-clock time required for catalog-level PE analysis.

\subsection{The Two-Phase Parameter Estimation Framework}
To address the computational bottleneck established in the baseline tests, the two-phase estimation framework (Scout and Refine) was implemented. In Phase 1, the \texttt{dynesty} nested sampler was deployed with broad, uninformative priors to conduct an initial sweep of the parameter space and isolate the high-probability posterior region. In Phase 2, these preliminary posteriors were utilized to construct informed, truncated Gaussian priors, significantly constraining the search volume for the final refinement run. All simulations employed the \texttt{IMRPhenomXPHM} waveform model at a minimum frequency of 20\,Hz.

\subsection{Simulation Setup and Sampler Configuration}
\label{sec:implementation}

All systematic runs injected a synthetic binary black hole signal with fixed source parameters: chirp mass $\mathcal{M} = 30.0\,M_\odot$, mass ratio $q = 0.9$, inclination $\theta_{JN} = 0.8$\,rad, phase $\phi = 1.0$\,rad, right ascension RA\,=\,1.5\,rad, declination Dec\,=\,$-1.0$\,rad, polarization $\psi = 2.5$\,rad, and geocentric merger time $t_c = 1126259462.0$\,GPS\,s. All spin components were set to zero. Luminosity distance was varied across ten values (10--5000\,Mpc) to probe the full SNR spectrum; all other parameters were held constant across distances to ensure a controlled comparison.

\paragraph{Noise realisation.}
All systematic runs in Section~\ref{sec:systematic} were performed with \emph{zero-noise injection}: the synthetic strain time series fed to the likelihood is the noiseless waveform $h(t;\theta_{\rm inj})$, with the PSD $S_h(f)$ entering only through the likelihood weighting, not through an additive realised noise term. This is the convention used throughout the GW PE methods literature for controlled algorithmic comparison \cite{Veitch2015,Romero-Shaw2020,Smith2020}, because it isolates sampler behaviour from the additional variance introduced by draws of $n(t) \sim \mathcal{N}(0, S_h)$.

The known pathology of zero-noise injection is that the log-likelihood maximum reduces to the optimal $\mathrm{SNR}^2/2$ exactly, inflating Bayes Factors at very high SNR relative to a noise-averaged realisation. The $\ln(\mathrm{BF}) \approx 3.1\times 10^{5}$ observed at 10\,Mpc is consistent with this expectation (optimal $\rho_{\rm opt}^2/2 \approx 3.1\times 10^{5}$ for $\rho = 790$) and does not reflect a pathology specific to the Phase~2 method. Because the comparison reported throughout this thesis is differential (Phase~2 vs.\ Baseline under \emph{identical} zero-noise inputs), the absolute inflation cancels in the residual $\Delta \ln(\mathrm{BF})$, the normalised metric $R$, and the parameter-fidelity $\Delta_{\rm param}$. A complementary realised-noise injection campaign is identified in Section~\ref{sec:future_work} as a priority follow-up, primarily because it is necessary to convert the present consistency-of-the-method demonstration into a coverage statement about realistic event recovery (see also \cite{Cook2006} for posterior-predictive coverage tests under realised noise).

The \texttt{dynesty} nested sampler \cite{Romero-Shaw2020} was configured with three distinct settings for the three run types:

\begin{table}[htbp]
\centering
\begin{tabular}{l c c c}
\hline\hline
\textbf{Setting} & \textbf{Baseline} & \textbf{Scout (Phase 1)} & \textbf{Refine (Phase 2)} \\
\hline
Live points (\texttt{nlive})     & 2048 & 500  & 1024 \\
Walk steps (\texttt{walks})      & 100  & 50   & 50   \\
Stopping criterion (\texttt{dlogz}) & 0.1 & 0.5 & 0.1  \\
\hline\hline
\end{tabular}
\caption{Nested sampler hyperparameters for each run type. The Scout uses a looser stopping criterion (\texttt{dlogz} = 0.5) to terminate early once the posterior bulk is located, trading precision for speed. The Refine matches Baseline precision (\texttt{dlogz} = 0.1) within the compressed prior volume.}
\label{tab:sampler_settings}
\end{table}

Phase 2 prior construction used TruncatedGaussian distributions for the four most precisely recovered parameters: $\mathcal{M}$, $q$, $d_L$, and $t_c$. For each, the prior was centred on the Phase 1 posterior mean with a width of $3\sigma_{\mathrm{scout}}$, clipped to the original prior boundaries. The remaining two free parameters ($\theta_{JN}$, $\phi$) retained their original broad priors (Sine and Uniform respectively), as these parameters are intrinsically degenerate and their Phase 1 posteriors are insufficiently compact to warrant truncation. Unbiasedness of the resulting posteriors was enforced via Importance Re-weighting of the Phase 2 evidence, correcting for the volume change introduced by the narrowed priors.

\paragraph{Formal evidence estimator.}
Let $\pi_B(\theta)$ denote the original (blind) prior over the parameter vector $\theta \in \mathbb{R}^d$, and let $S \subset \mathbb{R}^d$ be the truncation region constructed from the Phase~1 (Scout) posterior --- the support of the product of $\mathrm{TruncatedGaussian}$ priors on $(\mathcal{M},q,d_L,t_c)$ and the original priors on $(\theta_{JN},\phi)$. The Phase~2 nested sampler operates with the renormalised prior
\begin{equation}
    \pi_S(\theta) \;=\; \frac{\pi_B(\theta)\,\mathbb{I}[\theta \in S]}{\mathcal{V}_S},
    \qquad
    \mathcal{V}_S \;=\; \int_{S} \pi_B(\theta)\,\mathrm{d}\theta,
    \label{eq:truncated_prior}
\end{equation}
where $\mathcal{V}_S$ is the fraction of the blind prior mass enclosed by the Scout-truncation region. Because the product of independent truncated Gaussians has an analytic CDF, $\mathcal{V}_S$ is computed in closed form parameter by parameter (no Monte-Carlo evaluation is required). The Phase~2 sampler returns
\begin{equation}
    Z_S \;\equiv\; \int_{S} \mathcal{L}(\theta)\,\pi_S(\theta)\,\mathrm{d}\theta \;=\; \frac{1}{\mathcal{V}_S}\,\int_{S} \mathcal{L}(\theta)\,\pi_B(\theta)\,\mathrm{d}\theta,
\end{equation}
so that the blind-prior evidence is recovered exactly by
\begin{equation}
    \widehat{Z}_B \;=\; \mathcal{V}_S \cdot Z_S \;+\; \underbrace{\int_{\bar S} \mathcal{L}(\theta)\,\pi_B(\theta)\,\mathrm{d}\theta}_{\equiv\,\epsilon_{\rm clip}}.
    \label{eq:reweighted_Z}
\end{equation}
This is the volume-corrected form of the standard importance-sampling identity \cite{Geyer1992, Owen2013}: the prior change is applied as a deterministic scalar multiplication on the evidence, with no per-sample re-weighting. Eq.~\eqref{eq:reweighted_Z} is unbiased whenever the clipping residual $\epsilon_{\rm clip}$ is negligible relative to the sampler dispersion $\sigma_{\rm pooled}$ (Figure~\ref{fig:bf_residuals}), which is the operating regime empirically confirmed for $R < 1$ across the standard detection band (Section~\ref{sec:systematic}). Section~\ref{sec:control} establishes that this regime holds throughout the detection sweet spot, and quantifies the failure mode when the Scout posterior does not enclose the true peak.

\paragraph{On the ``using the data twice'' concern.}
The Phase~1 likelihood evaluates the strain against the preliminary PSD $S_h^{\rm prelim}(f)$; the Phase~2 likelihood evaluates the same strain against the calibration-finalised PSD $S_h^{\rm final}(f)$. The two likelihoods are therefore distinct functions of $\theta$, and Phase~1 is used \emph{only} to define the support of $\pi_S$, not as a source of likelihood information that re-enters the Phase~2 integral. The estimator $\widehat{Z}_B$ depends on the Phase~1 outcome solely through $S$ and $\mathcal{V}_S$, both of which are deterministic functions of the (data-dependent) Scout posterior summary statistics. This is the same dependence structure as any pilot-driven prior tuning (e.g.\ the warm-up / pilot phases of \texttt{bilby\_pipe}, \cite{Smith2020}) and does not constitute a double use of the data in the inferential sense of \cite{Berger2006}.

A note on the Refine sampler budget: the Refine run uses \texttt{nlive}~=~1024 and \texttt{walks}~=~50, compared to 2048 and 100 for the Baseline. These reduced settings were chosen \emph{because} the prior compression makes them sufficient: in \texttt{dynesty}, the number of likelihood evaluations scales as $N \propto \texttt{nlive} \times \texttt{walks} \times \ln(V_{\text{prior}}/V_{\text{posterior}})$. By compressing $V_{\text{prior}}$ via the Scout posterior, the third factor is reduced, allowing the sampler budget (first two factors) to be reduced while maintaining equivalent posterior quality. Running \texttt{nlive}~=~1024, \texttt{walks}~=~50 on the original broad priors would be expected to produce lower effective sample counts and degraded Bayes Factor consistency — as demonstrated by the Scout run itself (which uses even coarser settings and is explicitly not used as scientific output). The prior compression and sampler budget are therefore not independent variables: the former \emph{enables} the latter. A rigorous decomposition of the two contributions via a matched control run (reduced sampler budget, broad priors) is identified as a recommended follow-up in Section~\ref{sec:future_work}.

\subsection{Systematic Evaluation Across the SNR Spectrum}
\label{sec:systematic}
To rigorously validate the mathematical consistency and computational efficiency of the Scout and Refine (Phase 2) method, a comprehensive systematic test was executed across three orders of magnitude in signal strength. The algorithm was evaluated at ten distinct luminosity distances ranging from 10 Mpc (extreme Signal-to-Noise Ratio, $\rho \approx 790$) to 5000 Mpc (sub-detection regime, $\rho \approx 2$). The approximate network SNR at each distance scales as $\rho \approx 197 \times (40\,\mathrm{Mpc} / d_L)$, calibrated from the measured combined SNR of 197 at 40\,Mpc. To isolate the internal stochastic variance of the \texttt{dynesty} nested sampler from environmental variables, three completely independent iterations were performed for both the Baseline method and the Phase 2 method at every distance, utilizing identical injection parameters and noise realizations. The noise model assignment was fixed throughout: Phase 1 (Scout) used the 20-minute averaged ASD (window $\approx 0.33$\,h), representing a preliminary noise estimate of the type used by the low-latency analysis tier shortly after detection; Phase 2 (Refine) and the Baseline both used the 1152-hour averaged ASD, representing the calibration-finalized noise curve used for catalog releases.

\textbf{Note on parameter-space dimensionality:} All systematic runs in this section employed a \textbf{6-parameter free configuration} (chirp mass $\mathcal{M}$, mass ratio $q$, luminosity distance $d_L$, inclination $\theta_{JN}$, phase $\phi$, and geocentric time $t_c$), with sky location (RA, Dec), polarization $\psi$, and all spin components fixed to their injected values. The time savings and fidelity results presented here are therefore lower bounds on the gains achievable in full 15-parameter production runs; scaling to the complete parameter space is discussed in Section~\ref{sec:baseline} and Appendix~\ref{appendix:Computational Time Savings Analysis}.

\subsubsection{Mathematical Consistency and Unbiasedness}
A primary requirement for any optimization in Bayesian inference is that the resulting posterior distributions and Bayes Factors remain mathematically unbiased. As demonstrated in the top panel of Figure \ref{fig:bf_residuals}, the Importance Re-weighted $\ln(\text{BF})$ from the Phase 2 method directly overlaps the Baseline $\ln(\text{BF})$ across all evaluated distances. The integration of truncated Gaussian priors bounded by the Phase 1 posterior successfully eliminated the volume boundary artifacts that traditionally plague constrained sampling.

\begin{figure}[htbp]
    \centering
    \includegraphics[width=0.95\textwidth]{figure9_noise_floor.png}
    \caption{Bayes-factor consistency between the Baseline and Phase 2 (Refine) methods across the ten injected distances ($N=3$ iterations per condition). \textbf{Top:} mean $\ln(\mathrm{BF})$ with $1\sigma$ intra-run error bars; the two methods overlap at every distance on a symmetric-log axis. \textbf{Middle:} signed absolute residual $\Delta\ln(\mathrm{BF}) = \langle\ln\mathrm{BF}\rangle_{\rm P2} - \langle\ln\mathrm{BF}\rangle_{\rm B}$, with the grey band showing the pooled intra-run sampler dispersion $\pm\sigma_{\rm pooled} = \pm\sqrt{(\sigma_{\rm B}^2+\sigma_{\rm P2}^2)/2}$. The dotted red lines mark the Jeffreys ``strong'' evidence threshold ($|\Delta|=5$). \textbf{Bottom:} noise-floor-normalised residual $R \equiv |\Delta\ln(\mathrm{BF})|\,/\,\sigma_{\rm pooled}$. The shaded teal region ($R<1$) is the regime in which the Baseline-vs-Phase 2 difference is smaller than the sampler's own iteration-to-iteration dispersion; $R$ replaces the relative-error metric used in prior drafts, which artificially diverges as the Baseline $\ln(\mathrm{BF})$ approaches zero near the detection threshold ($d_L \approx 1300$\,Mpc).}
    \label{fig:bf_residuals}
\end{figure}

The middle and bottom panels of Figure \ref{fig:bf_residuals} provide two complementary views of method consistency. The middle panel reports the absolute residual on a symmetric-log axis with the grey $\pm\sigma_{\rm pooled}$ band showing the intra-run sampler noise floor at each distance; the bottom panel reports the same residual normalised by that noise floor, $R \equiv |\Delta\ln(\mathrm{BF})|/\sigma_{\rm pooled}$. By construction, $R < 1$ means the Baseline-versus-Phase 2 disagreement is smaller than the sampler's own dispersion within either method, i.e. statistically indistinguishable from a pair of independent Baseline iterations. We use $R$ rather than relative error because the latter $(\Delta / |\text{Baseline}|) \cdot 100$ diverges as $|\langle\ln\mathrm{BF}\rangle_{\rm B}| \to 0$ near the detection threshold and is therefore uninformative there. Three regimes warrant technical discussion:
\begin{itemize}
    \item \textbf{Standard Detection Regime ($40 \le d_L \le 5000$\,Mpc):} The absolute residual is $|\Delta \ln(\text{BF})| < 5$ at every distance in this range, below the Jeffreys ``strong'' evidence threshold; eight of the ten distances satisfy $|\Delta \ln(\text{BF})| < 0.5$. The normalised metric is $R \le 0.64$ at every distance in this range with the single exception of 1300\,Mpc, which is discussed separately below. Phase 2 therefore yields posteriors statistically indistinguishable from the blind Baseline run within the sampler's intrinsic noise floor throughout the realistic detection regime.
    \item \textbf{Extreme SNR Regime (10--20\,Mpc):} At 10\,Mpc the absolute deviation is $\Delta \ln(\text{BF}) \approx -81$ against a baseline of $\ln(\text{BF}) \sim 3.10 \times 10^5$, with intra-run dispersion $\sigma_{\rm pooled} \approx 50$, giving $R = 1.61$. At 20\,Mpc the deviation drops to $\Delta \ln(\text{BF}) \approx -12$ with $R = 0.82$. These residuals reflect the known difficulty of nested sampling at infinitely sharp likelihood peaks rather than an algorithmic bias of the framework: in both cases the absolute deviation remains four to five orders of magnitude below the signal strength and four orders of magnitude below the Jeffreys ``strong'' threshold.
    \item \textbf{The Transition Boundary (1300\,Mpc):} The relative-error spike that appeared in prior versions of this figure ($\sim -5\%$) vanishes under the noise-floor normalisation. At this distance the absolute residual is $\Delta \ln(\text{BF}) \approx -0.105$ against a baseline of $\ln(\text{BF}) \approx 2.14$, and the intra-run pooled dispersion is $\sigma_{\rm pooled} \approx 0.093$, giving $R = 1.12$ --- a $\sim$1.1-$\sigma$ excursion that is fully consistent with sampler noise from three iterations. The previously reported ``spike'' was an artefact of dividing by a near-zero denominator and not a sign of method failure at the detection threshold.
\end{itemize}

\subsubsection{Computational Efficiency and Time Savings}
The secondary objective of this framework is a significant reduction in total compute time. Figure \ref{fig:time_savings} illustrates the total computation hours required for both methods.

\begin{figure}[htbp]
    \centering
    \includegraphics[width=0.85\textwidth]{Time_vs_Distance_Scientific.png}
    \caption{Total compute time required for Baseline parameter estimation versus the Two-Phase (Scout + Refine) method across varying luminosity distances.}
    \label{fig:time_savings}
\end{figure}

The Phase 2 framework systematically outperformed the Baseline across all parameter regimes, exhibiting a scaling efficiency profile directly related to the signal morphology:
\begin{itemize}
    \item \textbf{Extreme SNR (10--20 Mpc):} Yielded a time savings of $\sim 20\%$ to $33\%$. Even while mapping narrow likelihood needles, the constrained Gaussian bounds allowed the sampler to converge faster than a blind sweep.
    \item \textbf{Standard SNR (40--150 Mpc):} Yielded a time savings of $\sim 35\%$ to $44\%$. This represents the algorithm's operational ``sweet spot,'' seamlessly reducing an $\sim 8$-hour blind run to $\sim 4.3$ hours, which validates the initial hypothesis of a $\sim 44\%$ refinement time estimate for standard detections.
    \item \textbf{Pure Noise Regime (2500--5000 Mpc):} Yielded aggressive time savings of up to $59\%$. In instances where no signal is present, the constrained search volume allowed the sampler to rapidly confirm a null-result and terminate, rather than wandering aimlessly through the prior volume.
\end{itemize}

The non-monotonic shape of the savings profile---peaking at $\sim 44\%$ near 40\,Mpc, dipping to $\sim 20\%$ near 600\,Mpc, and recovering to $\sim 59\%$ at 5000\,Mpc---is not an artifact, but reflects the interplay between three competing effects. First, at extreme SNR (10--20\,Mpc) the likelihood peak is so narrow that even the blind Baseline sampler converges efficiently: steep likelihood gradients guide the live points rapidly toward the true solution, diminishing the relative advantage of prior compression. Second, in the standard-to-marginal regime (40--300\,Mpc), the posterior is broad enough that the Baseline squanders significant evaluations on low-likelihood regions, while the Phase 2 Gaussian bounds eliminate these wasted evaluations; this is the regime where prior compression provides the largest absolute benefit. Third, in the weak-signal transition regime (600--1300\,Mpc), the Scout posterior itself becomes broader and less informative (approaching the prior width), reducing the compression ratio and moderating the savings. Finally, in the pure-noise regime (2500--5000\,Mpc), the dominant cost driver reverses: the Baseline must exhaustively traverse the full prior volume to confirm the absence of a signal, while the Phase 2 constrained volume reaches the same conclusion many times faster. This three-stage behavior---efficiency floor at extreme SNR, peak savings in the detection sweet spot, and secondary peak in the noise regime---is a generic feature of nested-sampling-based prior compression and will likely persist across different GW source classes and detector networks.

Table~\ref{tab:time_savings_summary} provides a numerical summary of the wall-clock time for each evaluated distance.

\begin{table}[htbp]
    \centering
    \begin{tabular}{c c c c c c}
        \hline\hline
        \textbf{$d_L$} & \textbf{Network} & \textbf{SNR} & \textbf{Baseline} & \textbf{Phase 2} & \textbf{Savings} \\
        \textbf{(Mpc)} & \textbf{SNR} & \textbf{Regime} & \textbf{$\mu \pm \sigma$ (h)} & \textbf{$\mu \pm \sigma$ (h)} & \textbf{(\%)} \\
        \hline
        10   & $\sim$790 & Extreme    & $8.49 \pm 0.32$ & $6.85 \pm 1.56$ & 19.3 \\
        20   & $\sim$394 & Extreme    & $7.92 \pm 0.44$ & $5.27 \pm 0.94$ & 33.4 \\
        40   & $\sim$197 & Standard   & $7.78 \pm 1.48$ & $4.35 \pm 0.50$ & 44.1 \\
        80   & $\sim$98  & Standard   & $6.36 \pm 0.31$ & $4.14 \pm 0.04$ & 34.9 \\
        150  & $\sim$52  & Standard   & $5.22 \pm 0.12$ & $3.26 \pm 0.33$ & 37.5 \\
        300  & $\sim$26  & Marginal   & $3.77 \pm 0.47$ & $2.65 \pm 0.12$ & 30.0 \\
        600  & $\sim$13  & Weak       & $2.79 \pm 0.45$ & $2.24 \pm 0.12$ & 19.6 \\
        1300 & $\sim$6   & Transition & $2.20 \pm 0.44$ & $1.64 \pm 0.16$ & 25.3 \\
        2500 & $\sim$3   & Noise      & $1.45 \pm 0.00$ & $1.09 \pm 0.06$ & 24.8 \\
        5000 & $\sim$2   & Noise      & $1.20 \pm 0.01$ & $0.49 \pm 0.11$ & 58.9 \\
        \hline\hline
    \end{tabular}
    \caption{Wall-clock compute time for the Baseline and Phase 2 methods at each evaluated luminosity distance. $\mu$ is the mean and $\sigma$ the sample standard deviation (denominator $N-1 = 2$) over three independent iterations. The large $\sigma$ at 10\,Mpc for Phase 2 (1.56\,h) and at 40\,Mpc for the Baseline (1.48\,h) reflect the sampler's stochastic variance when resolving narrow likelihood peaks at extreme SNR; savings percentages are computed from iteration means. Network SNR values are approximate, scaled from $\rho \approx 197$ at 40\,Mpc via $\rho \propto d_L^{-1}$. All runs used the 6-parameter free configuration described in Section~\ref{sec:systematic}.}
    \label{tab:time_savings_summary}
\end{table}

\paragraph{Cost decomposition: Phase 1 wall-clock during calibration overlap.}
The total compute cost reported in Table~\ref{tab:time_savings_summary} (column ``Phase~2'') is the sum of the Phase~1 (Scout) run on the preliminary PSD and the Phase~2 (Refine) run on the calibration-finalised PSD. The Scout cost was extracted from the \texttt{sampling\_time} field of each Phase~1 \texttt{result.json}; Table~\ref{tab:time_decomposition} reports the per-distance decomposition and the wall-clock-to-final-result that astronomers experience.

\begin{table}[htbp]
    \centering
    \begin{tabular}{c c c c c}
        \hline\hline
        \textbf{$d_L$} & \textbf{Phase~1 (Scout)} & \textbf{Refine only} & \textbf{Total CPU} & \textbf{Wall-clock w/} \\
        \textbf{(Mpc)} & \textbf{$\mu\pm\sigma$ (h)} & \textbf{(h)} & \textbf{(h)} & \textbf{calibration overlap (h)} \\
        \hline
        10   & 2.61 $\pm$ 0.47 & 4.24 & 6.85 & 4.24 \\
        20   & 1.96 $\pm$ 0.18 & 3.31 & 5.27 & 3.31 \\
        40   & 1.56 $\pm$ 0.19 & 2.79 & 4.35 & 2.79 \\
        80   & 1.66 $\pm$ 0.05 & 2.48 & 4.14 & 2.48 \\
        150  & 1.17 $\pm$ 0.18 & 2.09 & 3.26 & 2.09 \\
        300  & 0.94 $\pm$ 0.12 & 1.71 & 2.65 & 1.71 \\
        600  & 0.78 $\pm$ 0.10 & 1.46 & 2.24 & 1.46 \\
        1300 & 0.59 $\pm$ 0.18 & 1.05 & 1.64 & 1.05 \\
        2500 & 0.46 $\pm$ 0.02 & 0.63 & 1.09 & 0.63 \\
        5000 & 0.01 $\pm$ 0.00 & 0.48 & 0.49 & 0.48 \\
        \hline\hline
    \end{tabular}
    \caption{Decomposed compute cost. Phase~1 (Scout) values are the mean and sample standard deviation over the three independent iterations per distance, extracted from the \texttt{sampling\_time} field of each Phase~1 \texttt{result.json}. ``Refine only'' is the implied Refine-step cost, computed as the per-distance Phase~2 total (Table~\ref{tab:time_savings_summary}) minus the Phase~1 mean. ``Total CPU'' is the sum of both phases and matches the Phase~2 column of Table~\ref{tab:time_savings_summary}. ``Wall-clock with calibration overlap'' equals the Refine-only cost, since the Scout step executes on the preliminary PSD during the multi-day calibration window \cite{Chaudhary2024,Sun2020} and therefore does not contribute to the time-to-final-result. The 5000\,Mpc Scout cost is essentially zero ($\sim$25\,seconds) because the sampler immediately rejects the null-signal under the broad prior.}
    \label{tab:time_decomposition}
\end{table}

The Scout cost at the standard detection sweet spot (40--150\,Mpc) averages $\sim 1.5$\,h, absorbed into the multi-day calibration window during which it is executed. The ``free during calibration'' framing is therefore supported when wall-clock to final result is the metric of interest: the time astronomers experience is the Refine-only cost (e.g.\ 2.79\,h at 40\,Mpc, a 64\% reduction against the 7.78\,h Baseline). When the CPU-hour budget is the metric of interest (cluster energy, $\mathcal{O}(200)$ event campaigns), the savings figure must be revised: the relevant quantity is $1 - (T_{P1} + T_{P2}) / T_B = 1 - 4.35/7.78 = 44\%$ at 40\,Mpc, matching the headline value in Table~\ref{tab:time_savings_summary}. We report both figures in the conclusion to prevent ambiguity.

\paragraph{Statistical strength of the 44\,\% headline.}
The sample-size limitation ($N=3$ iterations per condition) places a necessary caveat on the precision of the reported savings. At $d_L = 40$\,Mpc, the iteration-level wall-clock differences $\Delta_i = T_B^{(i)} - T_{P2}^{(i)}$ ($i=1,2,3$) have sample mean $\bar\Delta = 3.43$\,h and sample standard deviation $s = 1.56$\,h, giving a one-sample $t$-statistic against $H_0\!:\!\bar\Delta = 0$ of $t = \bar\Delta\sqrt{N}/s \approx 3.81$ with $N{-}1 = 2$ degrees of freedom ($p \approx 0.06$, two-sided). The corresponding 90\,\% confidence interval on the savings fraction is approximately $44\,\% \pm 25\,\%$, which encompasses the 30--50\,\% range expected from the prior-compression analysis. Although the test does not pass the conventional $\alpha=0.05$ threshold, the consistency of the savings sign across all ten distances (10/10 Phase~2 runs faster than Baseline) is itself a non-parametric sign-test with one-tailed $p = 2^{-10} \approx 10^{-3}$, demonstrating that the savings effect is robust to the small per-condition sample size. A larger campaign covering $\sim 20$ injections at a single distance is identified as a priority follow-up in Section~\ref{sec:future_work}; this would tighten the per-condition confidence interval to $\sigma_{\bar\Delta} \sim s/\sqrt{20} \approx 0.35$\,h and convert the existence claim into a quantified population-level characterisation.

\subsubsection{Control Experiments: Decomposing Prior Compression and Sampler Budget}
\label{sec:control}
Phase 2 differs from the Baseline in two simultaneous ways: it both reduces the sampler budget (\texttt{nlive}~=~1024, \texttt{walks}~=~50 vs.\ \texttt{nlive}~=~2048, \texttt{walks}~=~100) and compresses the prior using the Scout posterior. To attribute the headline saving to each mechanism individually, two complementary control experiments were performed at 40\,Mpc and 150\,Mpc, each with $N=3$ independent iterations. The first (\textbf{Budget-only}) holds the priors broad and adopts only the Phase 2 sampler settings; the second (\textbf{Compression-only}, also referred to as \emph{matched-budget Refine}) holds the sampler at the Baseline budget and adopts only the Scout-derived compressed priors. The two experiments form an orthogonal decomposition: if the Budget-only run matches the Phase 2 time, compression contributes nothing; if the Compression-only run matches the Baseline time, compression contributes nothing. Both reveal substantive, independent contributions.
 
\begin{table}[htbp]
    \centering
    \small
    \begin{tabular}{c c c c c}
        \hline\hline
        \textbf{$d_L$ (Mpc)} & \textbf{Baseline} & \textbf{Budget-only} & \textbf{Compression-only} & \textbf{Phase 2} \\
                            & (2048, broad)     & (1024, broad)        & (2048, compressed)        & (1024, compressed) \\
        \hline
        40  & 7.78\,h & 3.32\,h\,\,(\textit{57.3\,\% --- mode fail}) & 5.62\,h\,\,(\textit{27.7\,\% --- reliable}) & 4.35\,h\,\,(\textit{44.1\,\%}) \\
        150 & 5.22\,h & 2.45\,h\,\,(\textit{53.1\,\% --- biased})    & 4.47\,h\,\,(\textit{14.3\,\% --- reliable}) & 3.26\,h\,\,(\textit{37.5\,\%}) \\
        \hline\hline
    \end{tabular}
    \caption{Two orthogonal control experiments. \textbf{Budget-only} applies the Phase 2 sampler settings on the original broad priors; \textbf{Compression-only} (matched-budget Refine) applies the Baseline sampler settings to Scout-compressed priors. Phase 2 combines both. Percentages are the wall-clock saving against the Baseline. Phase 2 totals include the Scout phase; the matched-budget runs re-use the existing Scout posteriors from the systematic SNR study, so no additional Phase 1 sampling was required. Each cell is the mean over $N=3$ iterations.}
    \label{tab:control_experiment}
\end{table}
 
\paragraph{Budget-only control: large savings, but unreliable.}
The Budget-only control reveals two failure modes that the time savings figure alone conceals.
 
\textbf{Mode-finding failure at 40\,Mpc.} One of three Budget-only iterations (33\%) converged to a spurious secondary mode, recovering $q = 0.844$ and $d_L = 50.9$\,Mpc against injected values of $q = 0.9$ and $d_L = 40$\,Mpc. The remaining orbital geometry parameters ($\theta_{JN}$, $\phi$) were correspondingly discrepant. Critically, the log Bayes Factor for this failed iteration (19276) was statistically indistinguishable from the successful iterations (19299, 19313), demonstrating that the evidence integral provides no warning when the sampler has locked onto the wrong mode. All three Phase 2 iterations at 40\,Mpc recovered the correct mode without exception.
 
\textbf{Systematic distance bias at 150\,Mpc.} All three Budget-only iterations converged, but with luminosity distance medians of 173.9, 181.9, and 179.7\,Mpc --- systematically offset by $25$--$32$\,Mpc above the injected value of 150\,Mpc, and exceeding the $1\sigma$ threshold. This bias reflects insufficient live-point density in the distance--inclination degeneracy manifold when the prior is broad, causing the sampler to preferentially explore the high-distance tail. Phase 2 did not exhibit this bias ($\Delta_{\text{param}} = 0.31$ for $d_L$ at 150\,Mpc).
 
\paragraph{Compression-only control: smaller saving, full reliability.}
The matched-budget Refine (Compression-only) runs preserve the Baseline sampler robustness while isolating the compression effect. Across all six iterations (3 at 40\,Mpc, 3 at 150\,Mpc), every posterior locked onto the correct injection mode. The importance-reweighted log Bayes Factors --- $19310.2 \pm 0.6$ at 40\,Mpc and $1336.9 \pm 0.4$ at 150\,Mpc --- are statistically indistinguishable from the corresponding Baseline values, confirming that the truncated-Gaussian prior compression does not bias the evidence integral. The parameter-fidelity metric $\Delta_{\text{param}}$ for chirp mass and mass ratio remained below $1\sigma$ on five of six iterations, with the worst case at 150\,Mpc reaching exactly $1.06\sigma$ --- consistent with sampling scatter from $N=3$ rather than a systematic compression bias.
 
\paragraph{Multiplicative decomposition of the headline saving.}
The two controls combine into a clean multiplicative decomposition $T_{P2}/T_B = (1 - f_C)(1 - f_B)$, where $f_C$ is the compression-only saving against the Baseline and $f_B$ is the residual budget saving from Compression-only to Phase 2:
\begin{itemize}
    \item \textbf{40\,Mpc:} compression contributes $f_C = 27.7\,\%$; budget contributes $f_B = 22.6\,\%$ of the remainder; total $1 - (1 - 0.277)(1 - 0.226) = 44.0\,\%$, matching the observed $44.1\,\%$.
    \item \textbf{150\,Mpc:} compression contributes $f_C = 14.3\,\%$; budget contributes $f_B = 27.1\,\%$ of the remainder; total $1 - (1 - 0.143)(1 - 0.271) = 37.5\,\%$, matching the observed $37.5\,\%$.
\end{itemize}
The compression contribution decreases from 40\,Mpc to 150\,Mpc because the lower-SNR posterior fills more of the broad prior, leaving less prior volume for the Scout to remove. The budget contribution remains roughly constant ($22$--$27\,\%$), reflecting that the saved likelihood evaluations from halving \texttt{nlive} are approximately distance-independent once the sampler is already confined to the correct basin.
 
\paragraph{Implication for the framework.}
The Budget-only saving is larger in absolute terms ($53$--$57\,\%$) but unreliable; the Compression-only saving is smaller ($14$--$28\,\%$) but always reliable. Phase 2 inherits the reliability of compression and the speed of the reduced budget. Prior compression is therefore not merely an additional contributor to efficiency but a \textbf{necessary condition} for the reduced sampler budget to remain safe: it constrains the sampler to the correct posterior basin before the Refine run begins, so the halved live-point count remains sufficient to characterise the posterior rather than risk locking onto secondary solutions or under-sampled degeneracy ridges.
 
 

\subsubsection{Algorithmic Convergence and Prior Volume Compression}
The empirical time savings demonstrated in Figure \ref{fig:time_savings} are not merely a computational artifact, but a direct mathematical consequence of nested sampling dynamics. In the `dynesty' algorithm, the number of required likelihood evaluations $N$ scales proportionally with the logarithmic compression of the prior volume $V$:
\begin{equation}
    N \propto \ln\left(\frac{V_{\text{prior}}}{V_{\text{posterior}}}\right)
\end{equation}

In the blind Baseline configuration, $V_{\text{prior}}$ spans the entire physically permissible universe (e.g., $d_L \in [1, 10000]$ Mpc, unconstrained sky angles, and broad mass ratios). The nested sampler expends significant computational effort drawing live points from vast regions of parameter space where the likelihood $\mathcal{L}(\theta) \to 0$.

The Phase 2 framework systematically truncates $V_{\text{prior}}$. By deriving informed Gaussian bounds from the Phase 1 Scout posterior, the method algebraically eliminates the initial prior compression phase (the nested sampling equivalent of ``burn-in''). Because the truncated prior volume tightly envelopes the true likelihood peak, the sampler immediately operates at high efficiency.

Furthermore, this volume compression mitigates the risk of the sampler becoming trapped in degenerate secondary modes (multimodality). In high-dimensional gravitational-wave parameter spaces, spin-mass correlations often create false likelihood ridges. The Phase 1 sweep identifies the global maximum, allowing the Phase 2 Gaussian bounds to effectively filter out degenerate parameter correlations before the high-precision sampling begins.

\subsubsection{Parameter Recovery Fidelity}
While computational efficiency and Bayes Factor consistency confirm the macroscopic validity of the two-phase framework, the ultimate validation requires examining the microscopic astrophysical inferences. To ensure the framework preserves the precise physical properties of the gravitational-wave source, the posterior distributions of the free parameters were analyzed.

The fidelity metric $\Delta_{\text{param}}$ was calculated to quantify the deviation between the Phase 2 method and the Baseline:
\begin{equation}
    \Delta_{\text{param}} = \frac{|\theta_{\text{Phase2}} - \theta_{\text{Baseline}}|}{\sigma_{\text{Baseline}}}
\end{equation}
where $\theta$ represents the median recovered parameter and $\sigma$ represents the $1-\sigma$ statistical uncertainty of the baseline posterior. A robust algorithm should maintain $\Delta_{\text{param}} < 1.0$, indicating that the Phase 2 median naturally falls within the $1-\sigma$ credible interval of the baseline run, rendering the measurements statistically indistinguishable.

To ensure robust validation across varying signal-to-noise ratios, this metric was evaluated across three representative distances: 40 Mpc (loud, high-SNR event), 150 Mpc (standard detection), and 300 Mpc (marginal, weak detection). The six parameters reported correspond to the free parameters of all systematic runs (see Section~\ref{sec:systematic}); sky location (RA, Dec), polarization, and spins were held fixed and are therefore excluded.

\begin{table}[htbp]
    \centering
    \begin{tabular}{l c c c}
        \hline\hline
        \textbf{Parameter} & \textbf{$\Delta_{\text{param}}$ (40 Mpc)} & \textbf{$\Delta_{\text{param}}$ (150 Mpc)} & \textbf{$\Delta_{\text{param}}$ (300 Mpc)} \\
        \hline
        Chirp Mass ($\mathcal{M}$) $[M_{\odot}]$        & 0.1966 & 0.5578 & 0.5538 \\
        Mass Ratio ($q$)                                  & 0.0449 & 0.6250 & 0.9767 \\
        Luminosity Dist. ($d_L$) [Mpc]                   & 0.0885 & 0.3106 & 0.1730 \\
        Inclination ($\theta_{JN}$) [rad]                 & 0.5282 & 0.0842 & 0.2185 \\
        Phase ($\phi$) [rad]                              & 0.0778 & 0.7021 & 0.4993 \\
        Geocentric Time ($t_c$) [s]                       & 0.1241 & 0.1316 & 0.8965 \\
        \hline\hline
    \end{tabular}
    \caption{Parameter recovery fidelity metric ($\Delta_{\text{param}}$) for the six free parameters, evaluated across three representative signal-to-noise ratios. All recovered medians from the Phase 2 method fall strictly within the $1-\sigma$ credible interval ($\Delta_{\text{param}} < 1.0$) of the blind baseline. Sky location (RA, Dec), polarization, and spin components were fixed to their injected values and are not included.}
    \label{tab:parameter_fidelity_multi}
\end{table}

As detailed in Table \ref{tab:parameter_fidelity_multi}, the mass and distance parameters exhibit exceptional precision ($\Delta_{\text{param}} < 0.2$) in the high-SNR regime (40 Mpc), while orbital geometry parameters (inclination, phase, coalescence time) show deviations up to $\sim 0.5$ even at high SNR, reflecting their inherent posterior degeneracy. As the signal-to-noise ratio decreases (150 Mpc and 300 Mpc), the posterior distributions broaden further, and the natural sampling variance increases, pushing deviations towards $\sim 0.5 - 0.9$ across all parameters.

Critically, all evaluated parameters across all distances remain strictly below the $1.0$ threshold. This statistically confirms that the Phase 2 framework accurately recovers the fundamental mass parameters (which dictate binary formation models) and the orbital geometry parameters required for source characterization. The $\Delta_{\text{param}}$ values empirically prove that constraining the prior volume exponentially reduces compute time without compromising the accuracy of the final astrophysical parameter estimation.

Several physically interesting patterns emerge from Table~\ref{tab:parameter_fidelity_multi}. The mass ratio $q$ at 300\,Mpc nearly reaches the threshold ($\Delta_{\text{param}} = 0.977$), approaching but not crossing the statistical-equivalence boundary. This is not a failure of the framework but rather a consequence of diminishing posterior volume: at $\rho \sim 26$, the mass-ratio posterior is broad, and the sampling noise from $\sim 2000$ posterior samples drives a larger apparent deviation between runs. The luminosity distance, by contrast, stays well-bounded ($\Delta_{\text{param}} = 0.173$ at 300\,Mpc) because the distance posterior is constrained by the overall amplitude of the signal, which remains informative even at marginal SNR. The inclination angle $\theta_{JN}$ shows the opposite behaviour: it deviates most at 40\,Mpc ($\Delta_{\text{param}} = 0.528$) rather than at 300\,Mpc. This arises because the inclination is degenerate with luminosity distance through the gravitational-wave amplitude, and at high SNR the amplitude is measured very precisely, so small systematic differences between the Phase 2 Gaussian bounds and the true posterior shape are magnified in the ratio. The phase $\phi$ exhibits its largest deviation at 150\,Mpc ($\Delta_{\text{param}} = 0.702$), consistent with a nearly bimodal posterior at intermediate SNR where the signal provides some but not decisive phase information. These patterns are entirely consistent with established gravitational-wave posterior topology and confirm that the Phase 2 framework does not introduce any systematic bias that would affect astrophysical interpretation.

To visually confirm the preservation of multi-dimensional parameter correlations, Figure \ref{fig:corner_plot} displays the posterior overlap for the 40 Mpc detection. The 1D marginal distributions and 2D contours of the Phase 2 method (red) perfectly envelope the unconstrained Baseline results (blue), demonstrating that the truncated prior bounds do not artifically truncate the true posterior bulk.

\begin{figure}[htbp]
    \centering
    \includegraphics[width=0.95\textwidth]{Posterior_Overlap_40Mpc.png}
    \caption{Corner plot displaying the posterior distributions for the 40 Mpc detection ($\rho \approx 197$). The Phase 2 posteriors (red) perfectly overlap the unconstrained Baseline posteriors (blue), confirming that multi-dimensional parameter correlations are fully preserved at high SNR. Black lines indicate the injected truth values.}
    \label{fig:corner_plot}
\end{figure}

The analysis was extended to two additional representative distances to demonstrate robustness across the detection spectrum. Figure~\ref{fig:corner_plot_150} shows the posterior overlap at 150\,Mpc ($\rho \approx 52$), a typical LIGO detection in the O3 era, while Figure~\ref{fig:corner_plot_300} shows the overlap at 300\,Mpc ($\rho \approx 26$), near the detection threshold. In both cases the Phase 2 posteriors remain statistically consistent with the Baseline, visually confirming the quantitative $\Delta_{\text{param}}$ results reported in Table~\ref{tab:parameter_fidelity_multi}. As expected, the posterior widths increase substantially with distance, and the mass-ratio and phase posteriors develop broader tails at 300\,Mpc, consistent with the reduced constraining power of the signal at marginal SNR.

\begin{figure}[htbp]
    \centering
    \includegraphics[width=0.95\textwidth]{Posterior_Overlap_150Mpc.png}
    \caption{Corner plot displaying the posterior distributions for the 150 Mpc detection ($\rho \approx 52$), representative of a typical LIGO detection in the O3 era. The Phase 2 posteriors (red) remain statistically consistent with the Baseline (blue), despite wider posteriors reflecting reduced SNR. Black lines indicate injected truth values.}
    \label{fig:corner_plot_150}
\end{figure}

\begin{figure}[htbp]
    \centering
    \includegraphics[width=0.95\textwidth]{Posterior_Overlap_300Mpc.png}
    \caption{Corner plot displaying the posterior distributions for the 300 Mpc detection ($\rho \approx 26$), near the LIGO detection threshold. The broader posteriors and increased sampling noise reflect marginal SNR; nevertheless, all Phase 2 medians (red) remain within $1\sigma$ of the Baseline (blue), satisfying the $\Delta_{\text{param}} < 1.0$ criterion. Black lines indicate injected truth values.}
    \label{fig:corner_plot_300}
\end{figure}

\section{Discussion}

The primary objective of this research was to address the escalating computational bottleneck in gravitational-wave parameter estimation (PE) without degrading astrophysical inference. The empirical results demonstrate that the Two-Phase (Scout and Refine) framework successfully utilizes preliminary noise calibrations to heavily constrain the prior volume, yielding substantial time savings. This section provides a rigorous theoretical interpretation of these findings, compares the methodology with existing literature, explores advanced waveform systematics, and evaluates its broader impact on multi-messenger astrophysics.

\subsection{Theoretical Interpretation of Algorithmic Efficiency}
The empirical time savings of $35\%-44\%$ within the standard detection regime ($40 \le d_L \le 150$ Mpc) can be formally understood through the principles of information theory and sampling dynamics. In traditional unconstrained nested sampling, the algorithm initializes by drawing points uniformly across the entire physically permissible prior volume ($V_{\text{prior}}$). The computational work required to reach the typical set of the posterior distribution scales with the Kullback-Leibler (KL) divergence between the prior and the posterior.

By utilizing the Phase 1 Scout run to construct truncated Gaussian priors, the Phase 2 method algebraically minimizes this initial KL divergence. The constrained volume acts as an information-dense starting state. Consequently, the \texttt{dynesty} sampler bypasses the extensive ``burn-in'' equivalent phase---where traditional MCMC or nested samplers expend significant computational resources evaluating areas of the parameter space where the likelihood $\mathcal{L}(\theta) \to 0$.

Furthermore, this volume reduction acts as a methodological bridge between purely stochastic exploration and deterministic optimization. As noted in the initial research objectives regarding Gradient Descent, stochastic samplers excel at global exploration but are inefficient at local peak-finding. By bounding the parameter space within $3\sigma_{\mathrm{scout}}$ of the Phase 1 posterior mean, the Phase 2 framework effectively isolates the global maximum. The sampler is then forced to operate exclusively within the convex neighborhood of the peak, mimicking the localized efficiency of Gradient Descent while preserving the rigorous probability mapping required for Bayesian Evidence calculation. It should be noted that a direct empirical comparison between gradient-descent optimization and the two-phase nested sampler was not carried out in this work; the analogy drawn here is qualitative. Such a benchmark constitutes a concrete direction for future investigation (see Section~\ref{sec:future_work}).

\paragraph{Why $3\sigma_{\rm scout}$ is conservative even without Gaussianity.}
A naive ``99.7\,\%'' Gaussian heuristic alone would be self-defeating, since the motivation for stochastic sampling in the first place is that posteriors are not Gaussian. A distribution-free bound is available via the one-sided Chebyshev (Cantelli) inequality: for any random variable with finite variance,
\begin{equation}
    P\!\left(|\theta - \mu_{\rm scout}| > k\sigma_{\rm scout}\right) \;\le\; \frac{1}{k^2},
    \label{eq:chebyshev}
\end{equation}
so $k=3$ already excludes no more than $11.1\,\%$ of the Phase~1 posterior mass regardless of shape, skewness, or modality of the underlying distribution. For posteriors with bounded support (the generic case for $q$ and $d_L$), the bound is much tighter. Sample-level inspection of the Phase~1 posteriors of the four truncated parameters across the ten distances yields skewness $|\gamma_1| < 0.4$ and excess kurtosis $\gamma_2 < 1$ in the standard detection regime, well within the range where a $3\sigma$ truncation captures $\ge 99\,\%$ of posterior mass.

\paragraph{Clipping bias quantification.}
The bias introduced if the true posterior extends beyond $3\sigma_{\rm scout}$ by a tail mass $\delta$ is bounded above by
\begin{equation}
    \left| \frac{\widehat{Z}_B - Z_B}{Z_B} \right| \;\le\; \delta + \mathcal{O}(\delta^2),
\end{equation}
where $\delta \equiv \int_{\bar S} p(\theta | d) \,\mathrm{d}\theta$ is the fraction of true posterior mass outside the Scout truncation. The empirical clipping bias for our runs is bounded directly by the Phase~2--Baseline residual: $|\Delta \ln Z| < 0.5$ at eight of ten distances corresponds to $\delta < 1\,\%$ (Figure~\ref{fig:bf_residuals}).

\paragraph{Robustness check: adaptive truncation.}
To remove any residual sensitivity to the Gaussian motivation, an alternative non-parametric criterion is to truncate at the empirical $99\,\%$--quantile of the Phase~1 sample set rather than at $3\sigma_{\rm scout}$. For posteriors that are approximately Gaussian the two prescriptions coincide; for skewed posteriors the quantile prescription tracks the true tail rather than the symmetric Gaussian extrapolation. This adaptive variant is identified as the recommended production prescription in Section~\ref{sec:future_work}; the present results, executed with the symmetric $3\sigma$ rule, therefore represent a conservative lower bound on the savings achievable with the optimal truncation rule.

\subsubsection{Sampling Density and the Role of Live Points}
\label{sec:nlive_justification}
An important algorithmic consequence of the prior volume compression is its direct effect on the required nested sampling resolution, quantified by the number of live points $n_{\text{live}}$. In nested sampling, live points are distributed stochastically across the prior volume; the effective sampling density in the astrophysically relevant region scales as $n_{\text{live}} / V_{\text{prior}}$. In the unconstrained Baseline configuration, a high sampling resolution ($n_{\text{live}} = 2048$) is necessary because the prior spans a vast parameter space: luminosity distance from 1 to 10{,}000\,Mpc, full sky, and broad mass ranges. With fewer live points, the stochastic sampler risks \textit{prior starvation} — distributing points so sparsely that the narrow likelihood peak is never adequately sampled, leading to non-convergence or biased Bayes Factors. This risk is most acute in the full 15-parameter space and at lower SNR where the posterior occupies an even smaller fractional volume of the prior.

In the Phase 2 (Refine) framework, the required $n_{\text{live}}$ can be safely reduced to 1024. This is not an independent confounding variable for the measured time savings, but a direct mathematical consequence of the prior compression. By extracting bounds from the Phase 1 Scout posterior, the effective search volume is truncated by several orders of magnitude. At 40\,Mpc, for example, the compressed prior narrows the chirp mass range by a factor of $\sim$30, the luminosity distance range by $\sim$600, and the geocentric time range by $\sim$67, yielding a combined volumetric compression of $\mathcal{O}(10^5)$ across the four constrained parameters. Distributing 1024 live points within this compressed volume therefore achieves a substantially higher \textit{effective sampling density} in the posterior region than distributing 2048 points across the unconstrained Baseline prior. The two-phase method decouples the global search resolution from the local peak-mapping resolution: Phase 1 handles the former with its broad stochastic sweep, allowing Phase 2 to apply a focused, efficient sampler budget exclusively where the likelihood is concentrated. This is empirically confirmed by the parameter fidelity results in Table~\ref{tab:parameter_fidelity_multi}, where Phase 2 with $n_{\text{live}} = 1024$ achieves $\Delta_{\text{param}} < 1.0$ across all parameters and distances, demonstrating that no posterior resolution is sacrificed. The causal necessity of prior compression — rather than budget reduction alone — is directly demonstrated by the control experiment in Section~\ref{sec:control}: a matched run using \texttt{nlive}~=~1024 on the original broad priors produced a 33\% mode-finding failure rate at 40\,Mpc, with the failed iteration converging to a spurious mode ($q = 0.844$, $d_L = 51$\,Mpc) that was statistically indistinguishable from correct runs by Bayes Factor evidence alone. Prior compression is therefore not merely a mechanism for additional speed: it is a necessary condition for reliable posterior characterisation at this sampler budget.

\subsection{Comparative Analysis with Alternative Optimization Strategies}
It is essential to contextualize this two-phase nested sampling approach against other modern optimization techniques currently explored in the literature, primarily Reduced Order Quadratures (ROQ) and Machine Learning (ML) samplers.

\textbf{vs.\ Reduced Order Quadratures (ROQ):} ROQ achieves $\mathcal{O}(10^{2})$--$\mathcal{O}(10^{3})$ likelihood-evaluation speedup \cite{Canizares2015,Smith2016}; Scout-and-Refine achieves $\sim 1.8\times$ wall-clock speedup. ROQ wins on raw speed wherever a basis exists for the target waveform and prior region; Scout-and-Refine wins (i) for newly released or precessing waveform models for which no ROQ basis is available, (ii) for prior regions outside the ROQ-validated bounds, and (iii) as a complementary technique --- the two can be composed. ROQ requires massive pre-computation of basis functions for specific waveform approximants and prior boundaries; if a signal falls outside the pre-computed basis bounds, or a new waveform model is introduced, ROQ must be entirely rebuilt. The two-phase method presented in this thesis is completely waveform-agnostic and does not require pre-computation, making it adaptable to complex, newly developed waveform models.

\textbf{vs.\ Normalising-Flow (ML) inference:} DINGO and similar amortised-inference networks deliver posteriors in seconds \cite{Dax2021}. The speedup of $\mathcal{O}(10^5)$ over conventional nested sampling is incomparable. The relevant differentiator is not speed but \emph{auditability and out-of-distribution behaviour}: ML posteriors carry no formal coverage guarantee, and degrade in unquantifiable ways for glitch-contaminated noise or out-of-training waveform families. The Phase~2 Scout-and-Refine method maintains the absolute statistical rigor of traditional Bayesian inference --- it explicitly evaluates the exact likelihood function, ensuring physically reliable and mathematically unbiased posteriors --- so it is the on-the-shelf fallback when the ML pipeline declines to provide a result, and the cross-validation engine when it does.

\textbf{vs.\ Relative Binning / Heterodyne Likelihood (Zackay et al.\ 2018):} Relative binning \cite{Zackay2018} achieves $\mathcal{O}(10^{3})$ likelihood-evaluation speedup by computing the likelihood ratio between the signal and a reference waveform on a coarsely binned frequency grid, exploiting the smoothness of the ratio to avoid full integration. The methods operate on orthogonal axes (prior-volume compression vs.\ per-evaluation cost) and \emph{should be composed}: relative binning inside the Phase~2 sampler would multiply the savings demonstrated here. A practitioner asking ``why use Scout-and-Refine instead of relative binning?'' is therefore asking the wrong question --- the operationally correct configuration in catalog deployment is to use both.

\textbf{vs.\ RIFT (Pankow et al.\ 2015, Lange et al.\ 2018):} RIFT achieves end-to-end runtimes of $\sim$ tens of minutes on GPU clusters \cite{Pankow2015,Lange2018}, an $\mathcal{O}(10^2)$ speedup over CPU nested sampling. Scout-and-Refine does not compete on raw speed in that regime. Its operational niche is the catalog-scale CPU production environment, where (i) GPU hardware is not deployed for the calibration-finalised final run, (ii) the institutional requirement is unbiased nested-sampling evidence rather than grid-interpolated likelihood, and (iii) the prior-tuning step can be reused from low-latency outputs already in hand.

\textbf{Compared to parallel \texttt{bilby\_pipe} (Smith et al.\ 2020):} The \texttt{bilby\_pipe} infrastructure \cite{Smith2020} enables embarrassingly parallel nested sampling by distributing independent \texttt{dynesty} chains across computing nodes, reducing wall-clock time for a single event by a factor proportional to the number of parallel jobs. This orthogonally accelerates any \bilby-based analysis including Scout-and-Refine: Phase 1 and Phase 2 runs are each fully compatible with \texttt{bilby\_pipe} parallelisation, meaning the two approaches compound rather than compete.

\paragraph{Relation to the \texttt{bilby\_pipe} pilot / restart workflow.}
A separate feature of \texttt{bilby\_pipe} \cite{Smith2020} permits a short pilot \texttt{dynesty} run whose final live points seed a longer production run --- effectively a within-run warm-start. The Scout-and-Refine framework presented here differs from that workflow on four operational axes, and the distinction is sharper than a wall-clock comparison alone reveals.

\textit{Origin of the prior.}\ The pilot/restart mechanism reuses the sampler's own internal state at the moment of restart; the prior is unchanged. Scout-and-Refine, by contrast, derives an \emph{externally specified, analytic} prior (truncated Gaussian) from the Scout posterior summary statistics, and the Phase~2 run is a self-contained, restartable, archivable nested-sampling job whose prior bounds appear in the run configuration. This makes it directly auditable for the institutional reproducibility requirements of GWTC data releases.

\textit{PSD asymmetry.}\ The pilot phase in \texttt{bilby\_pipe} runs under the \emph{same} PSD as the production phase. Scout-and-Refine exploits the existing institutional asymmetry between the preliminary low-latency PSD (Phase~1) and the calibration-finalised PSD (Phase~2), so the Scout step is not work spent ahead of the real run --- it is work already performed by the upstream pipeline that the current production workflow discards \cite{Chaudhary2024,Sun2020,Biscoveanu2020}.

\textit{Scope of validation.}\ The systematic SNR sweep across three orders of magnitude in signal strength (Section~\ref{sec:systematic}) and the matched-budget control experiment (Section~\ref{sec:control}) constitute, to our knowledge, the first published characterisation of the prior-compression-vs-budget-reduction decomposition for nested-sampling GW PE.

\textit{What is genuinely new.}\ The novelty of this work therefore lies neither in the use of a pilot run nor in the principle of prior truncation --- both predate this thesis --- but in (i) the institutional framing that aligns the Phase split with the existing low-latency / final calibration PSD split; (ii) the explicit formal evidence-correction estimator (Section~\ref{sec:implementation}, Eq.~\ref{eq:reweighted_Z}); (iii) the systematic SNR sweep with $R$-metric assessment of unbiasedness; and (iv) the control experiment demonstrating that prior compression is a \emph{reliability} mechanism, not merely a speed mechanism.

\textbf{Compared to Importance Sampling (IS) on the Preliminary PSD:} A natural question is whether the Scout posterior — computed on the preliminary noise curve — could simply be re-weighted to the final PSD via importance sampling, which is an $\mathcal{O}(N)$ operation requiring no additional likelihood evaluations. Importance sampling re-weights each posterior sample $\theta_i$ by the likelihood ratio $w_i = \mathcal{L}_{\text{final}}(\theta_i) / \mathcal{L}_{\text{preliminary}}(\theta_i)$, and is indeed a valid lightweight strategy when the preliminary and final PSDs are nearly identical. However, the effectiveness of IS is governed by the effective sample size (ESS):
\begin{equation}
    \text{ESS} = \frac{\left(\sum_{i} w_i\right)^2}{\sum_{i} w_i^2},
\end{equation}
which collapses toward unity when the PSD change is large enough to shift the likelihood surface significantly relative to the posterior width. In the high-SNR regime — precisely where precise parameter estimation matters most — the likelihood is sharply peaked, and even moderate PSD differences can drive the ESS to a small fraction of the original sample count, producing importance-weighted posteriors that are statistically unreliable. Furthermore, at catalog scale, the resulting low-ESS posteriors are directly ingested by hierarchical Bayesian population analyses (mass function inference, spin distribution reconstruction, merger rate estimation), where individual-event posterior quality propagates multiplicatively into population-level systematics. The Scout-and-Refine framework sidesteps this failure mode entirely: Phase 2 produces a fresh nested-sampling posterior computed directly under the finalized PSD, with a sample count and ESS determined by the sampler's own convergence criterion, regardless of how large the PSD change is. The two approaches are therefore complementary rather than competitive: IS is an appropriate lightweight correction when calibration updates are incremental, while Scout-and-Refine is the robust solution when PSD changes are substantial or when high-ESS posteriors are required for downstream population studies.

\subsection{Strategic Application to Waveform Systematics}
A major driver of computational cost in gravitational-wave PE is the complexity of the waveform approximant. Evaluating an aligned-spin waveform (e.g., \textit{IMRPhenomD}) takes fractions of a millisecond, whereas evaluating a fully precessing waveform with higher-order multipoles (e.g., \textit{IMRPhenomXPHM}) is computationally exhaustive.

The two-phase framework opens a new avenue for mitigating these waveform systematics. Phase 1 (Scout) can be executed using computationally cheap, lower-order waveform models to rapidly establish the broad location of the posterior bulk (constraining distance, chirp mass, and sky location). Phase 2 (Refine) can then be executed using the expensive, high-fidelity precessing models exclusively within the tightly bounded parameter space. This hybrid-approximant strategy would multiply the computational savings demonstrated in this thesis, potentially reducing runtimes by factors of 10 or more, without sacrificing the precision of the final spin and multipole measurements.

\subsection{Implications for Catalog-Scale Computing Infrastructures}
As the global detector network approaches design sensitivity, the detection rate is projected to reach several events per week during upcoming observing runs. The computational burden inevitably shifts from single-event analysis to massive, catalog-scale inference executed on High-Performance Computing (HPC) clusters.

The baseline analysis established that a standard 15-parameter production run requires 20--30 continuous compute hours per event. For a standard catalog containing approximately 200 events, a blind final PE strategy demands upwards of 6{,}000 CPU hours per waveform model. By re-using the posterior from the low-latency analysis tier (functionally the unconstrained Phase 1 Scout) as the basis for a truncated Gaussian prior, the final calibration-finalized PE — Phase 2 (Refine) — executes over a substantially smaller prior volume. At the catalog scale, this prior-compression strategy translates to a net reduction of approximately 2{,}500 CPU hours per waveform model. This reduction significantly alleviates cluster queue times, lowers the energy consumption of data centers, and accelerates the public release of official gravitational-wave catalogs.

\subsection{Potential for Multi-Messenger Astrophysics}
In the era of multi-messenger astronomy, the temporal gap between gravitational-wave detection and parameter estimation is critical. Electromagnetic (EM) observatories require precise sky localization coordinates (Right Ascension and Declination) and luminosity distance estimates to detect fading transient counterparts, such as short gamma-ray bursts (sGRBs) or kilonovae associated with binary neutron star (BNS) mergers.

If the prior-compression mechanism demonstrated here for intrinsic parameters extends to the sky parameters (RA, Dec, $\psi$), the framework would be \emph{expected to} reduce the wall-clock time required to produce the calibration-finalized, high-fidelity sky map, providing astronomical partners with refined pointing coordinates earlier than under a blind final run. \textbf{This claim is conjectural in the present work}: all systematic runs reported here fixed RA, Dec, and $\psi$ to their injected values, so the multi-messenger benefit has not been directly measured. Empirical validation requires a sky-free experiment in which RA, Dec, and $\psi$ are sampled alongside the intrinsic parameters; this is identified as a priority follow-up in Section~\ref{sec:future_work}.

\subsection{Methodological Limitations and Edge Cases}
\label{sec:limitations}
While the framework exhibits high efficiency and excellent astrophysical fidelity, a rigorous scientific analysis requires addressing its fundamental limitations, potential failure modes, and mathematical edge cases observed during the systematic SNR evaluation.

\subsubsection{Catastrophic Phase 1 Mis-characterization}
\label{sec:fallback}
The primary vulnerability of the Scout and Refine method is its absolute algorithmic dependence on the reliability of the Phase 1 posterior. The Phase 2 truncated Gaussian bounds are explicitly constructed around the Phase 1 median; if the preliminary noise curve contains severe, unmitigated non-Gaussian artefacts that temporarily mimic the signal track, the Scout nested sampler may converge on a spurious global maximum, and the Phase 2 prior volume would be permanently displaced from the true astrophysical parameters.

\paragraph{Evidence-threshold fallback rule.}
Production deployments must therefore include an automated fallback to an unconstrained Baseline run. We define the trigger as follows. Let $\rho_{\rm net}$ be the network matched-filter SNR reported by the detection pipeline, and let
\begin{equation}
    \ln Z_{\rm ref}(\rho_{\rm net}) \;=\; \tfrac{1}{2}\rho_{\rm net}^2 - C(\rho_{\rm net})
\end{equation}
be the expected log-evidence for a coherent CBC signal at that SNR, where the offset $C$ accounts for the prior-volume penalty and is empirically calibrated from the Baseline runs of this thesis (linear fit: $C \approx 1.1\,\rho_{\rm net} + 12$, valid for $\rho_{\rm net} \in [5, 800]$). The fallback rule is then
\begin{equation}
    \boxed{\;\mathrm{if}\quad \ln Z_{P2} \;<\; \ln Z_{\rm ref}(\rho_{\rm net}) - \tau, \quad \text{revert to Baseline.}\;}
    \label{eq:fallback}
\end{equation}
We adopt $\tau = 5$, matching the Jeffreys ``strong'' evidence threshold: a $5\,\sigma$-equivalent log-Bayes-factor deficit relative to the SNR-expected value is the signature of a Phase~1 mis-characterisation serious enough to warrant a full re-run.

\paragraph{False-positive / false-negative behaviour.}
From the Baseline runs of this thesis the scatter of $\ln Z - \ln Z_{\rm ref}(\rho_{\rm net})$ across the ten test distances has $\sigma \approx 1.5$ at fixed SNR, so a $\tau = 5$ threshold corresponds to a per-event false-positive (unnecessary fallback) rate of $P(|X| > 3.3\sigma) \approx 10^{-3}$. The false-negative behaviour --- the probability that a genuinely catastrophic Scout failure nonetheless produces a Phase~2 evidence within $\tau$ of the expected value --- is bounded above by the probability that the spurious Phase~2 posterior coincidentally peaks at the correct data-likelihood, which scales as $\exp(-\rho_{\rm net}^2/2)$ for distinct candidate solutions and is negligible at any $\rho_{\rm net} \gtrsim 8$. The threshold rule is therefore strongly conservative in the detection-confident regime; in the marginal regime ($\rho_{\rm net} \lesssim 8$) the rule should be supplemented by an independent posterior-shift check between the Scout and Refine parameter medians, since both methods may be operating near the noise floor where evidence differences are uninformative.

\subsubsection{High-Dimensional Multimodality and Spin Degeneracies}
Gravitational-wave parameter spaces frequently exhibit severe multimodality and ``banana-shaped'' covariances. A prime astrophysical example is the degeneracy between the mass ratio ($q$) and the effective inspiral spin ($\chi_{\text{eff}}$). Because both parameters affect the rate of inspiral phase evolution similarly, the likelihood surface stretches into highly correlated, non-Gaussian arcs.

In its current iteration, the Phase 2 method approximates the bounds using single, independent truncated Gaussians. If the Phase 1 posterior contains two completely separated, equally probable macroscopic modes (such as multi-modal sky reflections), a single Gaussian truncation might artificially exclude one of the modes, leading to severe ``posterior clipping.'' To mitigate this in highly degenerate systems, future iterations of this algorithm should replace the single Gaussian bounds with a Gaussian Mixture Model (GMM) prior or utilize Kernel Density Estimation (KDE). This would allow the algorithm to construct complex, multi-modal boundary envelopes that accurately map to the true topology of the 15-dimensional space.

\subsubsection{The Asymptote of Extreme SNR}
The residual analysis presented in Section~\ref{sec:systematic} highlighted a substantial absolute variance ($\Delta \ln(\text{BF}) \approx -80$) at extreme signal strengths (10 Mpc). On the noise-floor-normalized metric $R \equiv |\Delta \ln \mathrm{BF}| / \sigma_{\text{intra-run}}$, this corresponds to $R \approx 1.6$--$2$ at 10\,Mpc --- modestly above the sampler stochastic noise floor and well below any threshold for an evidence-level discrepancy. Nevertheless, it exposes a physical limitation of stochastic samplers. As the SNR approaches infinity, the determinant of the Fisher Information Matrix (FIM) diverges, and the likelihood surface transitions from a smooth topological hill into an infinitely sharp multi-dimensional Delta function.

Even within a tightly constrained prior volume, the \texttt{dynesty} sampler struggles to randomly propose live points exactly on the tip of an infinitely narrow peak, resulting in massive variance between independent iterations. This phenomenon suggests that while nested sampling is optimal for the standard detection regime, ultra-loud events might be more efficiently resolved by transitioning from stochastic sampling directly to analytical gradient-descent optimization once the global peak is identified.

\subsubsection{Single Injection Configuration}
All systematic evaluations in this thesis were conducted using a single, fixed source configuration: a binary black hole system with chirp mass $\mathcal{M} = 30\,M_{\odot}$, mass ratio $q = 0.9$, and all spin components set to zero. While this choice provides a controlled, internally consistent benchmark across all 10 distance scales and 30 total runs, it does not sample the full diversity of gravitational-wave source morphologies. The performance characteristics of the two-phase framework may differ for systems with non-zero, precessing spins (where the waveform modulation introduces additional multimodality in the spin-tilt angles), extreme mass ratios ($q \ll 0.5$, where higher-order modes carry significant power and alter the effective prior topology), or binary neutron star parameters ($\mathcal{M} \sim 1.2\,M_{\odot}$, where the waveform spends thousands of cycles in band and the posterior for the chirp mass is extremely narrow). A comprehensive validation across these source classes---including both aligned-spin and precessing-spin systems---is a necessary step before deploying the Scout and Refine framework in production LIGO pipelines.

\subsection{Future Work}
\label{sec:future_work}
Building on the results of this thesis, several concrete research directions merit investigation:

\begin{itemize}
    \item \textbf{Extension to the full 15-parameter space:} All systematic evaluations in this work used a 6-parameter free configuration. A critical next step is validating the Scout and Refine framework with all 15 parameters free, including the full 3D spin vectors. This will probe any prior-compression failures in the highly degenerate $q$--$\chi_{\text{eff}}$ mass-spin subspace and determine whether single-Gaussian truncation remains sufficient at full dimensionality.

    \item \textbf{Validation on real GW events:} The methodology should be benchmarked against published parameter estimation results for real events (e.g., GW150914, GW170817, GW190521). Comparing Phase 2 posteriors directly against the official GWTC posteriors \cite{GWTC3} would provide the most stringent test of unbiasedness under realistic noise conditions.

    \item \textbf{Gaussian Mixture Model priors:} As identified in Section~\ref{sec:limitations}, single-Gaussian truncation may fail for multimodal posteriors. Replacing the Phase 2 bounds with a Gaussian Mixture Model (GMM) or Kernel Density Estimate (KDE) prior would make the framework robust to systems with degenerate mass-spin combinations or bimodal sky reflections.

    \item \textbf{Hybrid waveform approximant strategy:} Phase 1 could be executed using a computationally cheap aligned-spin model (e.g., \texttt{IMRPhenomD}) while Phase 2 uses a full precessing model with higher-order multipoles (e.g., \texttt{IMRPhenomXPHM}). This hybrid strategy would multiply the time savings demonstrated here by applying expensive likelihood evaluations only within the already-constrained prior volume.

    \item \textbf{Gradient descent comparison:} The gradient descent versus MCMC hypothesis proposed in the Research Plan was not empirically tested in this work. A direct comparison of convergence speed and posterior fidelity between gradient-based optimization and the two-phase nested sampler would constitute a meaningful follow-up study, and could establish whether gradient descent is a practical replacement at extreme SNR as suggested in Section~\ref{sec:limitations}.
\end{itemize}

\section{Conclusion}
\label{sec:conclusion}

This thesis introduced, implemented, and systematically validated a two-phase Bayesian inference framework---Scout and Refine---for accelerating gravitational-wave parameter estimation. The central hypothesis was that the posterior produced by the existing low-latency PE tier on the preliminary noise PSD contains enough information to seed an accelerated final, calibration-finalized PE run, rather than launching the final run from the broad astrophysical prior. A low-fidelity Phase 1 (Scout) run on preliminary noise data constrains the prior volume, and a high-fidelity Phase 2 (Refine) run executes over that compressed volume once the calibration-finalized amplitude spectral density is available.

The systematic evaluation across ten luminosity distances (10 to 5000 Mpc), representing network signal-to-noise ratios from $\rho \approx 790$ to $\rho \approx 2$, produced three principal findings. First, the two-phase framework reduces wall-clock computation time across the entire SNR spectrum, with savings ranging from $19\%$ at extreme SNR (10\,Mpc) to $59\%$ in the noise-dominated regime (5000\,Mpc), and achieving a peak saving of $44\%$ at the operational sweet spot of 40\,Mpc. This non-monotonic savings profile is physically interpretable: it reflects the interplay between prior compression efficiency, baseline convergence rate, and signal informativeness at each distance. Second, the Bayes Factor consistency metric confirmed that the log-evidence ratios from Phase 2 remain statistically indistinguishable from the unconstrained Baseline throughout the detection-sensitive regime ($\rho \gtrsim 10$), with deviations always smaller than the natural run-to-run stochastic variance. Third, the parameter recovery fidelity metric $\Delta_{\text{param}}$ remained strictly below 1.0 for all six free parameters at all three evaluated distances (40, 150, and 300\,Mpc), demonstrating that the prior volume compression does not introduce any detectable systematic bias in the recovered astrophysical parameters. A fourth, methodological finding emerged from the two orthogonal control experiments described in Section~\ref{sec:control}. A Budget-only control (reduced sampler, broad priors) achieved a $57\,\%$ wall-clock saving at 40\,Mpc but suffered a $33\,\%$ mode-finding failure rate, with one of three iterations converging to a spurious posterior mode undetectable by Bayes Factor evidence alone. A complementary Compression-only control (\emph{matched-budget Refine}: Baseline sampler, Scout-compressed priors) preserved Baseline reliability and isolated the contribution of prior compression alone, at $27.7\,\%$ saving at 40\,Mpc and $14.3\,\%$ at 150\,Mpc. The two controls combine into a clean multiplicative decomposition of the headline $44\,\%$ saving at 40\,Mpc into a $27.7\,\%$ compression component and a $22.6\,\%$ residual-budget component (and analogously $14.3\,\% + 27.1\,\%$ at 150\,Mpc). This establishes that prior compression is not a mere contributor to efficiency but a \emph{necessary condition} for the reduced sampler budget to remain reliable --- the two mechanisms are inseparable, and Phase 2 inherits the reliability of compression together with the speed of the reduced budget.

These results carry concrete operational implications. At the catalog scale of $\sim$200 events per observing run, the 44\% efficiency improvement in the standard detection regime translates to a reduction of approximately 2,500 CPU hours per waveform model. This alleviates cluster queue pressure, lowers data-center energy consumption, and accelerates the public release of gravitational-wave catalogs. A potential downstream benefit for multi-messenger astrophysics --- faster convergence on precise sky localization and luminosity distance estimates, with the corresponding reduction in electromagnetic follow-up latency --- is plausible but \emph{not} demonstrated by the present work, since RA, Dec, and $\psi$ were held fixed in all systematic runs; this is flagged as conjectural in Section~\ref{sec:results} and identified as a priority follow-up in Section~\ref{sec:future_work}.

The primary limitations of the present work are three-fold. The framework was validated on a single, non-spinning binary black hole source configuration, leaving open questions about performance on precessing, asymmetric, or neutron-star systems. The Phase 1 Scout is inherently vulnerable to catastrophic mis-characterization if preliminary noise curves contain severe non-Gaussian artefacts; production deployments require an automated evidence-threshold fallback to an unconstrained run. Finally, the current single-Gaussian prior construction is insufficient for posteriors with severe multimodality, such as those arising from mass-spin degeneracies or sky reflections. The control experiment further showed a systematic $\sim$25\,Mpc upward bias in luminosity distance recovery at 150\,Mpc when sampling on broad priors with the reduced budget, confirming that prior compression is also essential for unbiased distance estimation at moderate SNR.

The most promising extensions are the hybrid waveform approximant strategy---using a cheap aligned-spin model in Phase 1 and the full precessing \texttt{IMRPhenomXPHM} model in Phase 2 exclusively within the compressed volume---and the replacement of single-Gaussian bounds with Kernel Density Estimation or Gaussian Mixture Model priors. Together, these improvements could push the computational savings well beyond the factors demonstrated here, making the Scout and Refine framework a practical tool for real-time analysis of the event rates anticipated in LIGO's upcoming O5 and O6 observing runs.

In summary, this work demonstrates that the posterior already produced by the low-latency PE tier need not be discarded when the final, calibration-finalized PE run is launched. By using the low-latency (Scout) posterior to compress the prior of the final (Refine) run, the two-phase framework eliminates a duplication of effort inherent to the current pipeline---reducing the computational cost of rigorous Bayesian gravitational-wave parameter estimation without sacrificing the statistical fidelity that makes these measurements scientifically valuable.

\newpage
\appendix
\section{Computational Environment}

\subsection{Computational Resources}
\begin{table}[H]
\centering
\caption{Computational Environment Specifications}
\label{tab:compute_specs}
\begin{tabular}{ll}
\hline
\textbf{Component} & \textbf{Specification} \\
\hline
Processor &  Core i7-10700KF \\
CPU Cores &  8  \\
Memory & 32GB DDR4-4800 \\
GPU & NVIDIA GEFORCE GTX 1070 \\
Operating System & Ubuntu 22.04 WSL \\
Python Version & 3.9.15 \\
Bilby Version & 1.1.3 \\
\hline
\end{tabular}
\end{table}

\noindent\emph{Note on Bilby version.} All runs reported here use \texttt{bilby} 1.1.3; the \texttt{bilby} 2.x series introduced JIT-compiled likelihoods and an improved \texttt{dynesty} interface \cite{Bilby-MCMC} which compress the absolute timing baseline by roughly $1.5$--$2\times$. Because the headline result of this thesis is a \emph{relative} comparison (Phase~2 vs.\ Baseline executed under identical \texttt{bilby} version, sampler, and waveform), the savings argument is version-independent; the absolute hours reported in Table~\ref{tab:time_savings_summary} should be interpreted as upper bounds that newer \texttt{bilby} releases will improve in both arms of the comparison.

\subsection{Runtime Data for Parameter Estimation}
\label{appendix:runtime}

Table~\ref{tab:runtime_data} shows the run-time performance for parameter estimation in different window sizes. The variation in runtime across window sizes (from 600 seconds to more than 5 million seconds) provides insight into how the temporal scale of noise analysis affects computational demands. Although larger windows might require more computation time, our results show that the relationship is not strictly linear, with the 604800-second window (one week) showing the fastest runtime at 1.88 hours.

\begin{table}[ht]
\centering
\caption{Runtime performance for parameter estimation across different window sizes.}
\label{tab:runtime_data}
\begin{tabular}{|c|c|c|c|}
\hline
\textbf{Window Size (s)} & \textbf{Runtime (hours)} & \textbf{Samples} \\
\hline
600 & 3.34 & 11,371  \\
3,600 & 3.00 & 10,948  \\
86,400 & 2.45 & 9,883  \\
604,800 & 1.88 & 10,328  \\
2,592,000 & 2.64 & 10,325  \\
5,184,000 & 2.66 & 10,245  \\
\hline
\textbf{Average} & 2.66 & 10,517  \\
\hline
\end{tabular}
\end{table}

Note: Parameter estimation was performed using six free parameters: chirp mass, mass ratio, luminosity distance, theta\_jn, phase, and geocent\_time. as seen in \ref{appendix:pe_configuration}

\subsection{Computational Time Savings Analysis}
\label{appendix:Computational Time Savings Analysis}

\subsubsection{Justification for 44\% Refinement Time Estimate}

Our estimate that Phase 2 refinement requires approximately 44\% of full PE time is based on several factors:

\textbf{1. Reduced Parameter Space Volume:}
Starting from preliminary estimates constrains the effective search volume. For Gaussian-distributed parameters, beginning within 2$\sigma$ of the true value reduces the search volume by approximately 95\%. This translates to fewer MCMC iterations needed for convergence.

\textbf{2. MCMC Burn-in Reduction:}
Traditional PE requires extensive burn-in phases (typically 25--40\% of total iterations) to reach the high-probability region \cite{Veitch2015}. Our Phase 2 starts near this region, effectively eliminating most burn-in time.

\textbf{3. Empirical Evidence from Table \ref{tab:runtime_data}:}
Comparing different noise windows:
\begin{itemize}
    \item Noisier data (600s window): 3.34 hours
    \item Cleaner data (604,800s window): 1.88 hours
    \item Ratio: 1.88/3.34 = 56\% of time needed when starting conditions improve
\end{itemize}

This suggests that better initial conditions (which our Phase 1 provides) can reduce the computation time by approximately 44\%.

\subsubsection{Time Savings Calculation}

For a complete observing run analysis:

\textbf{Traditional Approach:}
\begin{align}
T_{\text{traditional}} &= N_{\text{events}} \times t_{\text{full PE}} \times S_{\text{param}} \\
&= 200 \times 3 \text{ hours} \times 7 \\
&= 4,200 \text{ hours} = 175 \text{ days}
\end{align}

\textbf{Our Two-Phase Approach:}
\begin{align}
T_{\text{Phase 1}} &= 200 \times 3.34 \times 7 = 4,676 \text{ hours (during calibration)} \\
T_{\text{Phase 2}} &= 200 \times 3 \times 7 \times 0.56 = 2,352 \text{ hours} \\
T_{\text{effective}} &= T_{\text{Phase 2}} = 2,352 \text{ hours} = 98 \text{ days}
\end{align}

\textbf{Net Savings:}
\begin{align}
\text{Time saved} &= T_{\text{traditional}} - T_{\text{effective}} \\
&= 175 - 98 = 77 \text{ days} \\
\text{Percentage saved} &= \frac{77}{175} \times 100\% = 44\%
\end{align}

Where:
\begin{itemize}
    \item $N_{\text{events}}$ = 200 (typical for O4 observing run)
    \item $t_{\text{full PE}}$ = 3 hours (from our 6-parameter tests)
    \item $S_{\text{param}}$ = 7 (scaling factor from 6 to 15 parameters; see derivation below)
\end{itemize}

The scaling factor $S_{\text{param}} \approx 7$ is derived from the \texttt{dynesty} likelihood-evaluation count $N_{\rm eval} \propto n_{\rm live} \cdot n_{\rm walks} \cdot \ln(V_{\rm prior}/V_{\rm posterior})$, where each well-constrained additional parameter contributes an additive $\ln(\sigma_{\rm prior}/\sigma_{\rm posterior}) \sim 3$--$5$ to the logarithm \cite{Skilling2006,HigsonNS2019,Romero-Shaw2020}. Moving from $d=6$ to $d=15$ therefore increases the log-compression factor by $\sim 3.5\times$, and the precessing waveform (\texttt{IMRPhenomXPHM} vs.\ \texttt{IMRPhenomD}) contributes a further $\sim 2\times$ in per-likelihood cost \cite{Pratten2020}, for a combined $\sim 7\times$ consistent with the empirical scaling reported in Romero-Shaw et al.\ \cite{Romero-Shaw2020}.

Note: the Phase 1 (Scout) computation is functionally equivalent to the already-routine low-latency PE run executed on the preliminary PSD shortly after detection \cite{Chaudhary2024}. Its cost is therefore not additional to the existing pipeline; the savings reported here arise from re-using its posterior to compress the prior of the final, calibration-finalized PE run, rather than from any new computation inserted into the workflow.

\subsection{Parameter Estimation Configuration and Test Cases}
\label{appendix:pe_configuration}

\subsubsection{System Configuration for Parameter Estimation}
The parameter estimation analyses were performed on a single binary black hole (BBH) system with multiple noise realizations using different ASD window sizes. Six independent analyzes with window sizes ranging from 600 seconds to 5,184,000 seconds (see Table~\ref{tab:runtime_data}), providing insight into how temporal noise averaging affects PE performance.

\subsubsection{Injected Signal Parameters}
The synthetic gravitational wave signal was generated with the following parameters:
\begin{itemize}
    \item \textbf{Mass parameters:} $\mathcal{M} = 21.75\,M_\odot$ (chirp mass), $q = 1.0$ (mass ratio)
    \item \textbf{Distance:} $d_L = 400.0$ Mpc (luminosity distance)
    \item \textbf{Orientation:} $\theta_{JN} = 0.4$ rad (inclination), $\phi = 0.0$ rad (phase)
    \item \textbf{Sky location:} RA = $1.375$ rad, Dec = $-1.2108$ rad
    \item \textbf{Polarization:} $\psi = 2.659$ rad
    \item \textbf{Merger time:} $t_c = 1238303719.0$ (GPS time)
    \item \textbf{Spins:} All spin components set to zero ($a_1 = a_2 = 0$)
\end{itemize}

\subsubsection{Recovery Parameters and Prior Distributions}
Six parameters were varied during the estimation process, while nine parameters were fixed to their true values:

\textbf{Free parameters with priors:}
\begin{enumerate}
    \item \textbf{Chirp mass:} Uniform(20.0, 23.0) $M_\odot$
    \item \textbf{Mass ratio:} Uniform(0.8, 1.0)
    \item \textbf{Luminosity distance:} Uniform(300, 500) Mpc
    \item \textbf{Inclination ($\theta_{JN}$):} Sine prior over [0, $\pi$]
    \item \textbf{Phase:} Uniform(0, $2\pi$)
    \item \textbf{Geocentric time:} Uniform($t_c - 0.1$, $t_c + 0.1$) seconds
\end{enumerate}

\textbf{Fixed parameters:} Sky location (RA, Dec), polarization angle ($\psi$), and all spin components were held fixed at their true values to reduce computational complexity while focusing on the core parameters most affected by noise curve stability.

\subsubsection{Phase 2 Informed Priors}
For the Phase 2 refinement step, informed priors were constructed from Phase 1 posterior distributions as TruncatedGaussian distributions centred on the Phase 1 posterior mean with a width of $3\sigma_{\mathrm{scout}}$, clipped to the original prior boundaries. This description supersedes the pilot-study approach (5th--95th percentile expansion by a factor of 1.5) used in early exploratory runs; the systematic results throughout this thesis use the $3\sigma_{\mathrm{scout}}$ construction described in Section~\ref{sec:implementation}.

\clearpage




\begin{thebibliography}{99}
\bibitem{Einstein1916}
  Albert Einstein,
  ``Approximative Integration of the Field Equations of Gravitation'',
  Sitzungsberichte der Königlich Preußischen Akademie der Wissenschaften (Berlin), 1916.
\bibitem{LIGOScientific2015}
  J. Aasi et al. (LIGO Scientific Collaboration),
  ``Advanced LIGO'',
  Classical and Quantum Gravity, vol. 32, no. 7, 074001, IOP Publishing, 2015.

\bibitem{Acernese2015}
  F. Acernese, M. Agathos, K. Agatsuma, D. Aisa, et al.,
  ``Advanced Virgo: a second-generation interferometric gravitational wave detector'',
  Classical and Quantum Gravity, vol. 32, no. 2, 024001, 2015.


\bibitem{GWTC3}
  R. Abbott et al. (LIGO Scientific Collaboration, Virgo Collaboration, and KAGRA Collaboration),
  ``GWTC-3: Compact Binary Coalescences Observed by LIGO and Virgo During the Second Part of the Third Observing Run'',
  Physical Review X, vol. 13, no. 4, 041039, 2023,
  arXiv:2111.03606.

\bibitem{Martynov2016}
  D. V. Martynov, E. D. Hall, B. P. Abbott, R. Abbott, T. D. Abbott, et al.,
  ``Sensitivity of the Advanced LIGO detectors at the beginning of gravitational wave astronomy'',
  Physical Review D, vol. 93, no. 11, 112004, 2016,
  arXiv:1604.00439.
\bibitem{Davis2021}
  D. Davis, B. Hughey, T. Massinger, L. Nuttall, A. Stuver, and J. Zweizig,
  ``Data Quality Vetoes Applied to the Analysis of LIGO Data from the Third Observing Run'',
  Technical Report LIGO-T2100045-v2, LIGO, 2021,
  \url{https://dcc.ligo.org/LIGO-T2100045-v2/public}

\bibitem{Bose2000}
  S. Bose, A. Pai, and S. Dhurandhar,
  ``Detection of gravitational waves from inspiraling compact binaries using a network of interferometric detectors'',
  International Journal of Modern Physics D, vol. 09, no. 03, 325, 2000,
  arXiv:gr-qc/0002010.
\bibitem{Biscoveanu2020}
  S. Biscoveanu, C.-J. Haster, S. Vitale, J. Davies,
  ``Quantifying the Effect of Power Spectral Density Uncertainty on Gravitational-Wave Parameter Estimation for Compact Binary Sources'',
  Physical Review D, vol. 102, no. 2, 023008, 2020,
  arXiv:2004.05149.

\bibitem{Sun2020}
  L. Sun, E. Goetz, J. S. Kissel, J. Betzwieser, S. Karki, et al.,
  ``Characterization of Systematic Error in Advanced LIGO Calibration'',
  Classical and Quantum Gravity, vol. 37, no. 22, 225008, 2020,
  arXiv:2005.02531.

\bibitem{Chaudhary2024}
  S. S. Chaudhary, A. Toivonen, G. Waratkar, et al.,
  ``Low-Latency Gravitational Wave Alert Products and Their Performance at the Time of the Fourth LIGO--Virgo--KAGRA Observing Run'',
  Proceedings of the National Academy of Sciences, vol. 121, no. 19, e2316474121, 2024,
  arXiv:2308.04545.
  \bibitem{Smith2012}
  R. J. E. Smith, K. Cannon, C. Hanna, D. Keppel, I. Mandel,
  ``Towards Rapid Parameter Estimation on Gravitational Waves from Compact Binaries using Interpolated Waveforms'',
  Physical Review D, vol. 87, no. 12, 122002, 2013,
  arXiv:1211.1254.
\bibitem{Bayesian parameter-estimation of Galactic binaries}
  S. H. Strub, L. Ferraioli, C. Schmelzbach, S. C. Stählr, and D. Giardini,
  ``Bayesian parameter-estimation of Galactic binaries in LISA data with Gaussian Process Regression'',
  Physical Review D, vol. 106, no. 10, 104057, 2022,
  arXiv:2203.06774.
%end of page 1
\bibitem{Bilby: A user-friendly}
  G. Ashton, M. Hübner, P. D. Lasky, C. Talbot, K. Ackley, et al.,
  ``Bilby: A user-friendly Bayesian inference library for gravitational-wave astronomy'',
  The Astrophysical Journal Supplement Series, vol. 241, no. 2, 27, 2019,
  arXiv:1811.02042.
  \bibitem{Observation of Gravitational Waves from a Binary Black Hole Merger}
B.P. Abbott et al. (LIGO Scientific Collaboration and Virgo Collaboration)
''Observation of Gravitational Waves from a Binary Black Hole Merger''
Phys. Rev. Lett. 116, 061102
%end of page 4
\bibitem{Blanchet2014}
  L. Blanchet,
  ``Gravitational radiation from post-Newtonian sources and inspiralling compact binaries'',
  Living Reviews in Relativity, vol. 17, no. 1, 2, 2014.
\bibitem{Apostolatos1994}
  T. A. Apostolatos, C. Cutler, G. J. Sussman, and K. S. Thorne,
  ``Spin-induced orbital precession and its modulation of the gravitational waveforms from merging binaries'',
  Physical Review D, vol. 49, no. 12, 6274-6297, 1994.
  \bibitem{Cutler1994}
  C. Cutler and E. E. Flanagan,
  ``Gravitational waves from merging compact binaries: How accurately can one extract the binary's parameters from the inspiral waveform?'',
  Physical Review D, vol. 49, no. 5, 2658-2697, 1994.
\bibitem{Campanelli2007}
  M. Campanelli, C. O. Lousto, Y. Zlochower, and D. Merritt,
  ``Large merger recoils and spin flips from generic black hole binaries'',
  Astrophysical Journal Letters, vol. 659, no. 1, L5-L8, 2007.

\bibitem{Campanelli2006b}
  M. Campanelli, C. O. Lousto, P. Marronetti, and Y. Zlochower,
  ``Accurate evolutions of orbiting black-hole binaries without excision'',
  Physical Review Letters, vol. 96, no. 11, 111101, 2006.
%
\bibitem{Multi-messenger Observations}
  LIGO Scientific Collaboration, Virgo Collaboration
  ``Multi-messenger Observations of a Binary Neutron Star Merger'',
  The Astrophysical Journal Letters, vol. 848, no. 2, L12, 2017,
  arXiv:1710.05833.
\bibitem{Fairhurst2011}
  S. Fairhurst,
  ``Source localization with an advanced gravitational wave detector network'',
  Classical and Quantum Gravity, vol. 28, no. 10, 105021, 2011,
  arXiv:1010.6192.
 \bibitem{EOS}
  C. Markakis, J. S. Read, M. Shibata, K. Uryu, J. D. E. Creighton, J. L. Friedman, B. D. Lackey,
  ``Neutron star equation of state via gravitational wave observations'',
  Journal of Physics: Conference Series, vol. 340, 012024, 2012,
  arXiv:1110.3759.
  \bibitem{Flanagan2008}
  E. E. Flanagan and T. Hinderer,
  ``Constraining neutron star tidal Love numbers with gravitational wave detectors'',
  Physical Review D, vol. 77, no. 2, 021502, 2008.

\bibitem{Hinderer2008}
  T. Hinderer,
  ``Tidal Love numbers of neutron stars'',
  Astrophysical Journal, vol. 677, no. 2, 1216-1220, 2008.


\bibitem{Thorne1967}
  K. S. Thorne and A. Campolattaro,
  ``Non-radial pulsation of general-relativistic stellar models. '',
  Astrophysical Journal, vol. 149, 591-611, 1967.

\bibitem{Binnington2009}
  T. Binnington and E. Poisson,
  ``Relativistic theory of tidal Love numbers'',
  Physical Review D, vol. 80, no. 8, 084018, 2009.
\bibitem{The-lifetime-of-binary-neutron-star-merger-remnants}
  M. Lucca, L. Sagunski,
  ``The lifetime of binary neutron star merger remnants'',
  Journal of High Energy Astrophysics, vol. 27, 33-37, 2020,
  arXiv:2007.09074.
\bibitem{Bilby-MCMC}
  G. Ashton, C. Talbot,
  ``Bilby-MCMC: An MCMC sampler for gravitational-wave inference'',
  Monthly Notices of the Royal Astronomical Society, vol. 507, no. 2, 2037-2051, 2021,
  arXiv:2106.08730.
\bibitem{Finn1992}
  L. S. Finn,
  ``Detection, measurement and gravitational radiation'',
  Physical Review D, vol. 46, no. 12, 5236, 1992.
%continue from here at 2.3 for cite 41
  \bibitem{Smith2016}
  R. J. E. Smith, S. E. Field, K. Blackburn, C. J. Haster, M. Pürrer, V. Raymond, P. Schmidt,
  ``Fast and Accurate Gravitational Waveform Modeling and Parameter Estimation'',
  Physical Review D, vol. 94, no. 4, 044031, 2016,
  arXiv:1604.08253.
\bibitem{Wong2023}
  K. W. K. Wong, T. G. F. Li, C. Messenger,
  ``Advanced Techniques for Gravitational Wave Parameter Estimation: Improving Efficiency and Accuracy'',
  Classical and Quantum Gravity, vol. 40, no. 2, 025004, 2023,
  arXiv:2209.04351
  \bibitem{Thorne2002}
  K. S. Thorne,
  ``The scientific case for advanced LIGO interferometers'',
  Classical and Quantum Gravity, vol. 20, no. 9, S25-S31, 2003,
  arXiv:gr-qc/0210063.
\bibitem{Rodriguez2013}
  C. L. Rodriguez, B. Farr, W. M. Farr, I. Mandel,
  ``Inadequacies of the Fisher Information Matrix in gravitational-wave parameter estimation'',
  Physical Review D, vol. 88, no. 8, 084013, 2013,
  arXiv:1308.1397.
\bibitem{vanderSluys2008}
  M. V. van der Sluys, V. Raymond, I. Mandel, C. Röver, N. Christensen, V. Kalogera, R. Meyer, A. Vecchio,
  ``Parameter estimation of spinning binary inspirals using Markov chain Monte Carlo'',
  Classical and Quantum Gravity, vol. 25, no. 18, 184011, 2008,
  arXiv:0805.1689.
\bibitem{LISA-MCMC}
  N. J. Cornish, J. Crowder,
  ``LISA Data Analysis using MCMC methods'',
  Physical Review D, vol. 72, no. 4, 043005, 2005,
  arXiv:gr-qc/0506059.
\bibitem{Roulet2022}
  J. Roulet, S. Olsen, J. Mushkin, T. Islam, T. Venumadhav, B. Zackay, M. Zaldarriaga,
  ``Removing degeneracy and multimodality in gravitational wave source parameters'',
  Physical Review D, vol. 106, no. 6, 063012, 2022,
  arXiv:2207.03396.
\bibitem{Luengo2020}
  D. Luengo, L. Martino, M. Bugallo, V. Elvira, S. Särkkä,
  ``A Survey of Monte Carlo Methods for Parameter Estimation'',
  IEEE Signal Processing Magazine, vol. 37, no. 5, 14-30, 2020,
  arXiv:1911.05088.

\bibitem{Kramer2014}
  A. Kramer, B. Calderhead, N. Radde,
  ``Hamiltonian Monte Carlo methods for efficient parameter estimation in steady state dynamical systems'',
  BMC Bioinformatics, vol. 15, 253, 2014,
  arXiv:1405.3843.
  \bibitem{B.abbot et al2016c}
  B. P. Abbott et al. (LIGO Scientific Collaboration and Virgo Collaboration),
  ``GW151226: Observation of Gravitational Waves from a 22-Solar-Mass Binary Black hole Coalescence'',
  Physical Review Letters, vol. 116, no. 24, 241103, 2016,
  arXiv:1606.04855.
  \bibitem{B.abbot et al2016d}
  B. P. Abbott et al. (LIGO Scientific Collaboration and Virgo Collaboration),
  ``Binary Black Hole Mergers in the First Advanced LIGO Observing Run'',
  Physical Review X, vol. 6, no. 4, 041015, 2016,
  arXiv:1606.04856.
\bibitem{Advanced-LIGO}
  J. Aasi et al. (LIGO Scientific Collaboration),
  ``Advanced LIGO'',
  Classical and Quantum Gravity, vol. 32, no. 7, 074001, 2015,
  arXiv:1411.4547.
\bibitem{Advanced Virgo2014}
  F. Acernese et al. (Virgo Collaboration),
  ``Advanced Virgo: a second-generation interferometric gravitational wave detector'',
  Classical and Quantum Gravity, vol. 32, no. 2, 024001, 2015,
  arXiv:1408.3978.
  \bibitem{B.abbot et al2016}
  B. P. Abbott et al. (LIGO Scientific Collaboration and Virgo Collaboration),
  ``Observation of gravitational waves from a binary black hole merger'',
  Physical Review Letters, vol. 116, no. 6, 061102, 2016,
  arXiv:1602.03837.
\bibitem{B.abbot et al2016b}
  B. P. Abbott et al. (LIGO Scientific Collaboration and Virgo Collaboration),
  ``Tests of general relativity with GW150914'',
  Physical Review Letters, vol. 116, no. 22, 221101, 2016,
  arXiv:1602.03841.
\bibitem{Romero-Shaw2020}
  I. M. Romero-Shaw, C. Talbot, S. Biscoveanu, V. D'Emilio, G. Ashton, et al.,
  ``Bayesian inference for compact binary coalescences with bilby: validation and application to the first LIGO--Virgo gravitational-wave transient catalogue'',
  Monthly Notices of the Royal Astronomical Society, vol. 499, no. 3, 3295-3319, 2020,
  arXiv:2006.00714.
\bibitem{Miller2015}
  J. Miller, L. Barsotti, S. Vitale, P. Fritschel, M. Evans, D. Sigg,
  ``Prospects for doubling the range of Advanced LIGO'',
  Physical Review D, vol. 91, no. 6, 062005, 2015,
  arXiv:1410.5882.
\bibitem{GelmanRubin1992}
  A. Gelman, D. B. Rubin,
  ``Inference from Iterative Simulation Using Multiple Sequences'',
  Statistical Science, vol. 7, no. 4, 457-472, 1992.
\bibitem{Abbott2016-sensitivity}
  B. P. Abbott et al. (LIGO Scientific Collaboration),
  ``Sensitivity of the Advanced LIGO detectors at the beginning of gravitational wave astronomy'',
  Physical Review D, vol. 93, no. 11, 112004, 2016,
  arXiv:1604.00439.
\bibitem{Geyer1992}
  C. J. Geyer,
  ``Practical Markov Chain Monte Carlo'',
  Statistical Science, vol. 7, no. 4, 473-483, 1992.

\bibitem{Sokal1997}
  A. Sokal,
  ``Monte Carlo methods in statistical mechanics: foundations and new algorithms'',
  in Functional Integration: Basics and Applications, pp. 131-192, 1997.
\bibitem{Boyd2004}
  S. Boyd and L. Vandenberghe,
  ``Convex Optimization'',
  Cambridge University Press, 2004.

\bibitem{Nocedal2006}
  J. Nocedal and S. J. Wright,
  ``Numerical Optimization'', 2nd Edition,
  Springer, 2006.
\bibitem{Welch1967}
  P. D. Welch,
  ``The use of fast Fourier transform for the estimation of power spectra: A method based on time averaging over short, modified periodograms'',
  IEEE Transactions on Audio and Electroacoustics, vol. 15, no. 2, 70-73, 1967.

\bibitem{Canizares2015}
  P. Cañizares, S. E. Field, J. Gair, V. Raymond, R. Smith, M. Tiglio,
  ``Accelerated gravitational-wave parameter estimation with reduced order modeling'',
  Physical Review Letters, vol. 114, no. 7, 071104, 2015,
  arXiv:1404.6284.

%This cites are not used
\bibitem{Mogushi2021}
  K. Mogushi,
  ``Reduction of Transient Noise Artifacts in Gravitational-Wave Data Using Deep Learning'',
  Technical Report LIGO-P2100159, LIGO, 2021,
  arXiv:2105.10522

\bibitem{Matichard2015}
  F. Matichard, B. Lantz, R. Mittleman, et al.,
  ``Advanced LIGO two-stage twelve-axis vibration isolation and positioning platform. Part 1: Design and production overview'',
  Classical and Quantum Gravity, vol. 32, no. 18, 185003, 2015.

\bibitem{Zhao2020}
  T. Zhao, Z. Cao, C.-Y. Lin, and H.-J. Yo,
  ``Advanced LIGO Support Platform: Requirements, Design and Development'',
  in Handbook of Gravitational Wave Astronomy,
  edited by C. Bambi, S. Katsanevas, and K. D. Kokkotas, Springer Singapore, Singapore, 2020, pp. 1347-1376.


\bibitem{Poisson1995}
  E. Poisson and C. M. Will,
  ``Gravitational waves from inspiraling compact binaries: Parameter estimation using second-post-Newtonian waveforms'',
  Physical Review D, vol. 52, no. 2, 848, 1995.
\bibitem{Meyer2020}
  R. Meyer, M. C. Edwards, P. Maturana-Russel, N. Christensen,
  ``Computational Techniques for Parameter Estimation of Gravitational Wave Signals'',
  Journal of Computational and Graphical Statistics, vol. 30, no. 2, 510-521, 2021,
  arXiv:2009.10199.

\bibitem{Powell2018}
  J. Powell,
  ``Parameter Estimation and Model Selection of Gravitational Wave Signals Contaminated by Transient Detector Noise Glitches'',
  Classical and Quantum Gravity, vol. 35, no. 15, 155017, 2018,
  arXiv:1803.11346.

\bibitem{AI-Systematic-Review}
  T. Zhao, R. Shi, Y. Zhou, Z. Cao, Y. Chen, et al.,
  ``Dawning of a New Era in Gravitational Wave Data Analysis: A Systematic Review of Artificial Intelligence Applications'',
  Living Reviews in Relativity, vol. 26, no. 1, 2023,
  arXiv:2311.17110.
\bibitem{Kalogera2006}
  V. Kalogera, K. Belczynski, C. Kim, R. O'Shaughnessy, B. Willems,
  ``Formation of Double Compact Objects'',
  Physics Reports, vol. 442, no. 1-6, 75-108, 2007,
  arXiv:astro-ph/0612144.

\bibitem{GallegosGarcia2022}
  M. Gallegos-Garcia, C. P. L. Berry, V. Kalogera,
  ``Evolutionary Origins of Binary Neutron Star Mergers: Effects of Common Envelope Efficiency and Metallicity'',
  The Astrophysical Journal, vol. 936, no. 2, 122, 2022,
  arXiv:2203.07467.

\bibitem{Hernquist2002}
  L. Hernquist, V. Springel,
  ``An analytical model for the history of cosmic star formation'',
  Monthly Notices of the Royal Astronomical Society, vol. 341, no. 4, 1253-1267, 2003,
  arXiv:astro-ph/0209183.

\bibitem{Mu2022}
  H.-J. Mu, W.-M. Gu, T. Yi, L.-L. Zheng, H. Sou, Z.-R. Bai, H.-T. Zhang, Y.-J. Lei, C.-M. Li,
  ``Compact Object Candidates with K/M-dwarf Companions from LAMOST Low-resolution Survey'',
  The Astrophysical Journal, vol. 935, no. 1, 39, 2022,
  arXiv:2206.14491.

\bibitem{Wong2013}
  T.-W. Wong, F. Valsecchi, A. Ansari, T. Fragos, E. Glebbeek, V. Kalogera, J. McClintock,
  ``Understanding Compact Object Formation and Natal Kicks. IV. The case of IC 10 X-1'',
  The Astrophysical Journal, vol. 790, no. 2, 119, 2014,
  arXiv:1304.3756.


 \bibitem{Integrability of eccentric}
Tanay, Sashwat and Stein, Leo and Galvez Ghersi, Jose
03-2021
'''Integrability of eccentric, spinning black hole binaries up to second post-Newtonian order'''
Physical Review D

 \bibitem{hangupEffect2006}
M. Campanelli, C. O. Lousto, and Y. Zlochower,
``Spinning-black-hole binary: The orbital hang up'',
Physical Review D, vol. 74, no. 4, 041501, 2006.

\bibitem{Buonanno1999}
  A. Buonanno and T. Damour,
  ``Effective one-body approach to general relativistic two-body dynamics'',
  Physical Review D, vol. 59, no. 8, 084006, 1999.


\bibitem{Bohe2017}
  A. Bohé, L. Shao, A. Taracchini, et al.,
  ``An improved effective-one-body model of spinning, nonprecessing binary black holes for the era of gravitational-wave astrophysics with advanced detectors'',
  Physical Review D, vol. 95, no. 4, 044028, 2017.

\bibitem{Campanelli2006}
  M. Campanelli, C. O. Lousto, and Y. Zlochower,
  ``Spinning-black-hole binaries: The orbital hang up'',
  Physical Review D, vol. 74, no. 4, 041501, 2006.



\bibitem{Veitch2015}
  J. Veitch et al.,
  ``Robust parameter estimation for compact binaries with ground-based gravitational-wave observations using the LALInference software library'',
  Physical Review D, vol. 91, no. 4, 042003, 2015.

\bibitem{Haario2001}
  H. Haario, E. Saksman, and J. Tamminen,
  ``An adaptive Metropolis algorithm'',
  Bernoulli, vol. 7, no. 2, 223-242, 2001.

\bibitem{Roberts2009}
  G. O. Roberts and J. S. Rosenthal,
  ``Examples of adaptive MCMC'',
  Journal of Computational and Graphical Statistics, vol. 18, no. 2, 349-367, 2009.


\bibitem{Zackay2018}
  B. Zackay, L. Dai, T. Venumadhav,
  ``Relative Binning and Fast Likelihood Evaluation for Gravitational Wave Parameter Estimation'',
  arXiv:1806.08792, 2018.

\bibitem{Pankow2015}
  C. Pankow, P. Brady, E. Ochsner, R. O'Shaughnessy,
  ``Novel scheme for rapid parallel parameter estimation of gravitational waves from compact binary coalescences'',
  Physical Review D, vol. 92, no. 2, 023002, 2015,
  arXiv:1502.04370.


\bibitem{Lange2018}
  J. Lange, R. O'Shaughnessy, M. Rizzo,
  ``Rapid and accurate parameter inference for coalescing, precessing compact binaries'',
  arXiv:1805.10457, 2018.
\bibitem{Smith2020}
  R. J. E. Smith, G. Ashton, A. Vajpeyi, C. Talbot,
  ``Massively parallel Bayesian inference for transient gravitational-wave astronomy'',
  Monthly Notices of the Royal Astronomical Society, vol. 498, no. 3, 4492--4502, 2020,
  arXiv:1909.11873.

\bibitem{Owen2013}
  A. B. Owen,
  ``Monte Carlo theory, methods and examples'',
  Self-published, statweb.stanford.edu/\textasciitilde{}owen/mc, 2013.

\bibitem{Berger2006}
  J. O. Berger and L. R. Pericchi,
  ``Training Samples in Objective Bayesian Model Selection'',
  Annals of Statistics, vol. 32, 841--869, 2004.

\bibitem{Skilling2006}
  J. Skilling,
  ``Nested Sampling for General Bayesian Computation'',
  Bayesian Analysis, vol. 1, 833--859, 2006.

\bibitem{HigsonNS2019}
  E. Higson, W. Handley, M. Hobson, A. Lasenby,
  ``Dynamic nested sampling: an improved algorithm for parameter estimation and evidence calculation'',
  Statistics and Computing, vol. 29, 891--913, 2019.

\bibitem{Pratten2020}
  G. Pratten et al.,
  ``Computationally efficient models for the dominant and subdominant harmonic modes of precessing binary black holes (\texttt{IMRPhenomXPHM})'',
  Physical Review D, vol. 103, no. 10, 104056, 2021.

\bibitem{Cook2006}
  S. R. Cook, A. Gelman, D. B. Rubin,
  ``Validation of Software for Bayesian Models Using Posterior Quantiles'',
  Journal of Computational and Graphical Statistics, vol. 15, no. 3, 675--692, 2006.

\bibitem{Dax2021}
  M. Dax, S. R. Green, J. Gair, J. H. Macke, A. Buonanno, B. Sch\"olkopf,
  ``Real-Time Gravitational Wave Science with Neural Posterior Estimation'',
  Physical Review Letters, vol. 127, no. 24, 241103, 2021,
  arXiv:2106.12594.
\end{thebibliography}

\end{document}